% Este fichero es parte del Número 5 de la Revista Occam's Razor
% Revista Occam's Razor Número 5
%
% (c)  2010, The Occam's Razor Team
%
% Esta obra está bajo una licencia Reconocimiento 2.5 España de Creative
% Commons. Para ver una copia de esta licencia, visite 
% http://creativecommons.org/licenses/by/2.5/es/
% o envie una carta a Creative Commons, 171 Second Street, Suite 300, 
% San Francisco, California 94105, USA.
%
% Modificada para Revista GLUD Magazine (volumen 2 :: 2017)}
% ~~ Este es el archivo que se compila! ^^'


\documentclass[10pt, a4paper, twoside]{article}

% Paquetes... probablemente alguno no sea necesario
% SI INCLUYE ALGUNO POR FAVOR COLOQUE UN COMENTARIO AL FINAL DEL MISMO.

\usepackage[utf8]{inputenc}                                                   
\usepackage[greek,spanish]{babel}  % Símbolo del euro
\usepackage{graphicx}
\usepackage{a4, fancyhdr, multicol}
\usepackage{float}
\usepackage{pdftricks}
\usepackage{pstricks}
\usepackage{color}
\usepackage{pst-plot}
\usepackage{pst-eps}
\usepackage{wrapfig}
\usepackage{eso-pic}
\usepackage{textcomp} % Agregado para copyleft
\usepackage{multirow}
\usepackage{listings}
\usepackage{textpos}
\usepackage{epsf}
\usepackage{setspace}
\usepackage{hyperref}
\usepackage{colortbl}
\hypersetup{
    bookmarks=true,         % show bookmarks bar?
    colorlinks=true,        % false: boxed links; true: colored links
    linkcolor=red,          % color of internal links
    citecolor=green,        % color of links to bibliography
    filecolor=magenta,      % color of file links
    urlcolor=blue           % color of external links
}


\usepackage[T1]{fontenc} 


% Configuración de tamaño de página
\setlength{\parindent}{0in}
\setlength{\parskip}{0.1cm}
\setlength{\oddsidemargin}{0.05mm}
\setlength{\evensidemargin}{0.05mm}

\addtolength{\textwidth}{4cm}
\addtolength{\topmargin}{-3.5cm}
\addtolength{\textheight}{4.5cm}

\pagestyle{fancy}

% Configuración de Cabeceras Fancy Header
\fancyhead{}
\fancyfoot{} % clear all footer fields
\fancyfoot[LE]{\textbf{\textsf{\thepage|GLUD Magazine $\cdot$ Vol. 2, a\~no 2017 $\cdot$ Grupo GNU Linux Universidad Distrital}}}
\fancyfoot[RO]{\textbf{\textsf{GLUD Magazine $\cdot$ Vol. 2, a\~no 2017 $\cdot$ Grupo GNU Linux Universidad Distrital | \thepage }}}
\renewcommand{\footrulewidth}{0.4pt}

% Elimina l.neas de cabecera
% 
%\renewcommand{\headrulewidth}{0pt}

% Colores >.<
\definecolor{introcolor}{rgb}{0.8,1.0,0.8}
\definecolor{listcolor}{rgb}{0.2,0.4,0.2}
\definecolor{titlecolor}{rgb}{0.4,0.5,0.1}
\definecolor{excolor}{rgb}{0.8,0.8,0.8}
\definecolor{notecolor}{rgb}{0.4,0.4,1.0}

% ********************************************************
% Definición de Comandos y Entornos
% ********************************************************

% Comandos de uso general
% ---------------------------------------------------------
% Secciones títulos y subtítulos de cada página
\newcommand{\msection}[4]{
{\begin{flushright}{
{\psset{linecolor=black,linestyle=dotted}\psline(-17,0)}
\colorbox{#1}{
\begin{minipage}{#3\linewidth}
\center
  \textcolor{#2}{
    \textsf{\textbf{#4}}}
\end{minipage}
}}\end{flushright}}

\vspace{4mm}
}

\newcommand{\mtitle}[2]{
  {\resizebox{#1}{0.7cm}{\textbf{\textsf{#2}}}}
  \vspace{1mm}
}


\newcommand{\msubtitle}[2]{
  {\resizebox{#1}{0.5cm}{{\gray{\textbf{\textsf{#2}}}}}}
  \vspace{1mm}
}

% Principio de Página. Pone el cuadro superior con la sección
\newcommand{\bOpage}[3]{
  \msection{#1}{black}{#2}{#3}
  \begin{multicols}{2}
}

% Fin de página. Termina el entorno multicols
\newcommand{\eOpage}{
\pagebreak
\end{multicols}
}

% Fin e Inicio de Página. Sino utilizamos figuras fuera de las
% columnas del cuerpo principal, esta es la forma adecuada de marcar
% cada página
\newcommand{\ebOpage}[3]{
\eOpage
\bOpage{#1}{#2}{#3}
}

% Crea el cuadro de introducción al principio de cada artículo
\newcommand{\intro}[3]{
\colorbox{#1}{
  \begin{minipage}{.9\linewidth}
    \vspace{2mm}
    {{\resizebox{!}{1.0cm}{#2}}{#3}}
  \vspace{1mm}
  \end{minipage}
}
\vspace{4mm}
}

% Comando para introducir figura en entorno multicol
\newcommand{\myfig}[3]{
\begin{center}
  \includegraphics[width=#3\hsize,angle=#1]{#2}
  \nobreak
\end{center}}

% Caption para figuras en entorno multicol
\setcounter{figure}{1}
\newcommand{\mycaption}[1]{
  \begin{quote}
    {\small
    {{\sc Figura} \arabic{figure}: #1}
    }
  \end{quote}
  \stepcounter{figure}
}


% Comandos para utilizar tablas y figuras en entorno multicols
\makeatletter
\newenvironment{tablehere}
  {\def\@captype{table}}
  {}

\newenvironment{figurehere}
  {\def\@captype{figure}}
  {}
\makeatother

% Caption para figuras en entorno multicol sin contador
\newcommand{\nncaption}[1]{
  %\begin{quote}
    {\footnotesize{\textbf{
    {#1}
    }}}
  %\end{quote}
}

\newcommand{\sectiontext}[3]{\vspace{4mm}{{\textcolor{titlecolor}{\large{\textbf{\textsf{#3}}}}}}
\vspace{1mm}
}

\newcommand{\EOP}{\psframe[fillstyle=solid, fillcolor=titlecolor, linecolor=titlecolor](0,0)(4pt,4pt)}

% Entornos (begin... end)
% ----------------------------------------------------
% Entorno para introducir ejemplos
\newenvironment{mexample}{
  \vspace{2mm}
  \bgroup
  \tiny
}{
  \egroup
  \vspace{4mm}
}

% Entorno para introducir entradillas en el texto
\newenvironment{entradilla}{
  \vspace{5mm}
  \hrule 
  \vspace{2mm}
  \bgroup
  \LARGE
\begin{spacing}{0.6}
}{
\end{spacing}
  \egroup
  \vspace{5mm}
  \hrule
  \vspace{5mm}
}


% Entorno para introducir Biografia del Autor.
\newenvironment{biografia}[2]{
\vspace{10mm}
\hrule
\begin{minipage}[t]{3cm}
\myfig{0}{#1}{1} 
\end{minipage}
%\includegraphics[width=2.54 \hsize]{#1}
  \nobreak
\begin{minipage}[r]{5cm}
 \bgroup
  \small
\textbf{#2}
}{
  \egroup
  \vspace{3mm}
\end{minipage}
  \hrule
  \vspace{5mm}
}
% Entorno para introducir la bibliografia del articulo.
\newenvironment{bibliografia}{
\begin{thebibliography}{1}
}{
\end{thebibliography}
}



% **************************************************************************
% Comienza el documento
\begin{document}

% Portada, no utiliza Fancy Header e introduce imagen de portada con PStricks
%\pagestyle{empty}

%\rput(8,-14){{\resizebox{!}{30cm}{\epsfbox{images/portada_n5.eps}}}}

%\clearpage
%\pagebreak

% Imagen con el sumario en la siguiente página
%\rput(8,-14.0){\resizebox{!}{30cm}{{\epsfbox{images/sumario.eps}}}}

%\pagebreak

% Activa Fancy Headers stilo e incluye los distintos artículos
\pagestyle{fancy}

\definecolor{introcolor}{rgb}{0.6,0.7,0.3}

% Este fichero es parte del Número 5 de la Revista Occam's Razor
% Revista Occam's Razor Número 5
%
% (c)  2010 The Occam's Razor Team
%
% Esta obra está bajo una licencia Reconocimiento 2.5 España de Creative
% Commons. Para ver una copia de esta licencia, visite 
% http://creativecommons.org/licenses/by/2.5/es/
% o envie una carta a Creative Commons, 171 Second Street, Suite 300, 
% San Francisco, California 94105, USA.
% 
% Adaptada en el año 2017 para el Grupo GNU Linux Universidad Distrital
%

\rput(8,-1.6){\resizebox{!}{5cm}{{\epsfbox{images/general/editorial.eps}}}}
\rput(0.0,-13){\resizebox{7cm}{35.0cm}{{\epsfbox{images/general/bar.eps}}}}

\rput(0.7,-26.5){\resizebox{!}{0.9cm}{{\epsfbox{images/general/licencia.eps}}}}


\begin{textblock}{9.2}(3,0)
\begin{flushright}
{\resizebox{!}{1cm}{\textsc{Editorial}}}

\vspace{2mm}

Grupo GNU Linux \\Universidad Distrital\\Francisco José de Caldas
\end{flushright}
\end{textblock}

\vspace{4mm}


\definecolor{barcolor}{rgb}{0.9,0.9,0.9}

\begin{textblock}{30}(-1.5, -1)
\begin{minipage}{0.12\linewidth}
\sf\color{barcolor}
\begin{center}

\vspace{1cm}

\colorbox{black}{
{\resizebox{3cm}{0.7cm}{\textcolor{white}{\bf\sf\large GLUD}}}
}
{\resizebox{2.5cm}{0.4cm}{\textcolor{white}{\bf\sf\large Magazine}}}
%{\resizebox{2.5cm}{0.4cm}{\bf\sf\large Razor}}

\vspace{4mm}

{\bf Número 2.\\ Octubre 2017}

%\vspace{5cm}
\vspace{2cm}
\hrule

\vspace{7mm}

%{\bf Dirección: }

%\vspace{1mm}

%Colocar Dirección de la revista

%\vspace{2mm}

%{\bf Editores:}

%\vspace{1mm}

%Colocar comite editorial.

%\vspace{4mm}

{\bf Colaboradores:}

\vspace{1mm}

{\tt
Jose Noe Poveda\\
Diego Osorio\\
Steven Sierra\\
Andrés Acosta\\
Laydi Bautista\\
Fernando Pineda\\
Marlon Cárdenas\\
Sebastian Sanchez\\
Jesus David Romero\\
Alejandro Cortazar\\
Jorge Ulises Useche\\
Juan Carlos Velandía\\
Leidy Marcela Aldana}

\vspace{15.0mm}

%{\bf Maquetación y Grafismo}

%\vspace{1mm}

%Completar

\vspace{2mm}

\hrule

\vspace{4mm}

{\bf Publicidad}

\vspace{3mm}

Occam's Razor Direct

{\tt occams-razor@uvigo.es}

\vspace{4mm}

\hrule

\vspace{6mm}

{\bf Impresión}

Por ahora tu mismo\ldots Si te apetece

\vspace{3mm}

\hrule

\vspace{9mm}

\textcopyleft 2017 GLUD Magazine \\

Esta obra está bajo una licencia Reconocimiento 2.5 España de Creative
Commons. Para ver una copia de esta licencia, visite 

{\scriptsize \href{http://creativecommons.org/licenses/by/2.5/es/}{creativecommons.org}} 

\end{center}
\end{minipage}

\end{textblock}

\begin{textblock}{20}(3,2.0)

\begin{minipage}{.45\linewidth}
\colorbox{introcolor}{
\begin{minipage}{1\linewidth}

{{\resizebox{!}{1.0cm}{E}}{l Grupo GNU Linux Universidad Distrital agradece a cada uno de los lectores de esta revista, y a los desarrolladores del c\'odigo fuente de la misma, es decir, a los encargados de la revista Occam's Razor. 


\bigskip

}
}

\end{minipage}
}

\vspace{8mm}

El Grupo GNU Linux es un Grupo Lider en la apropiación, desarrollo, uso y difusión de tecnología, ciencia y cultura libre; el cual, hace parte de la Universidad Distrital Francisco José de Caldas.

\vspace{3mm}
GLUD Magazine fue fundada en el año 2011, seis años después, para este segundo volumen publica su segundo número para compartir artículos acerca de ciencia, tecnología  y cultura libre.

\bigskip

Como grupo, esperamos que sea agradable para usted leer esta revista y que sirva como aporte a la comunidad.

\medskip
%{\large Feliz 2017!}

\begin{flushright}
{\Large\sc{Grupo GLUD.}}
\end{flushright}

\vspace{5cm}

\bigskip


\end{minipage}

\bigskip

\colorbox{introcolor}{
\begin{minipage}{.45\linewidth}

\bigskip

{\footnotesize\sf{\color{white}
Las opiniones expresadas en los artículos, así como los contenidos de los mismos, son responsabilidad de los autores de éstos.

Puede obtener la versión electrónica de esta publicación, por medio de la página:\\
\begin{center}
	{\tt \href{https://glud.udistrital.edu.co/}{glud.udistrital.edu.co}}
\end{center}

Así como el {\em código fuente} de la misma y los distintos ficheros de datos asociados a cada artículo en el repositorio de la revista en Github:\\
\begin{center}
	{\tt \href{https://github.com/GLUD/GLUD_Magazine/}{github.com}}
\end{center}


}}

\bigskip

\end{minipage}
}


\end{textblock}

\pagebreak

%% Este fichero es parte del Número 5 de la Revista Occam's Razor
% Revista Occam's Razor Número 5
%
% (c)  2010 The Occam's Razor Team
%
% Esta obra está bajo una licencia Reconocimiento 2.5 España de Creative
% Commons. Para ver una copia de esta licencia, visite 
% http://creativecommons.org/licenses/by/2.5/es/
% o envie una carta a Creative Commons, 171 Second Street, Suite 300, 
% San Francisco, California 94105, USA.

% Este .tex contiene el contenido del articulo
% Seccion (Dejar en blanco)
%

% Incluye imagen del artículo (Debe ser diferente del Kokopelli! (°_°) Sin embargo, no es obligatoria esta imagen)

\rput(2.5,-2.3){\resizebox{!}{3.2cm}{{\includegraphics[]{images/nombre_articulo/glud.eps}}}}


%Imagen de el comienzo de el articulo, coordenadas desde %la parte superior izquierda del margen de la pagina
% Las imagenes quedan mejor en formato .eps, agradecería si alguien logra colocar imagenes en dicho formato y comparte la corrección.

% -------------------------------------------------
% Cabecera
\begin{flushright}

\msection{introcolor}{black}{0.35}{Seccion dejar igual}

\mtitle{10cm}{Titulo en Español del Artículo}

\msubtitle{8cm}{Titulo en Inglés del Artículo}

{\sf Nombre Del Autor(es), y profesión }

{\psset{linecolor=black,linestyle=dotted}\psline(-12,0)}
\end{flushright}

\vspace{2mm}
% -------------------------------------------------

\begin{multicols}{2}

\sectiontext{white}{black}{Resumen}

% Introducción
\intro{introcolor}{E}{n esta parte por favor ubicar el Resumen de su artículo, debe permitir identificar en forma rápida el contenido básico del trabajo; no debe ser mayor a 250 palabras. No debe aportar información o conclusión que no está presente en el texto, así como tampoco debe  citar referencias bibliográficas. Debe quedar claro el problema que se investiga y el objetivo del mismo. 
}

\sectiontext{white}{black}{Palabras clave: } Primera, Segunda, Tercera.


\sectiontext{white}{black}{I. Introducción}

Añadir la introducción del texto. \\
Gracias a todos los que comparten su trabajo!\\
\textbf{Cualquier comentario es bien recibido.}\\
Por si alguien necesita incluir código:


Código en Python.
\begin{verbatim}
for i in range (1, 5):
    print i
else:
    print "The for loop is over"
\end{verbatim}

Código en  C.

\begin{lstlisting}
#include <stdio.h>
int main(int argc, char* argv[]) {
  puts("Hola mundo!");
}
\end{lstlisting}

\vspace{2mm}
 
Código en Matlab (no debería incluir esto, pero venía con la plantilla original, quizá a alguien le sirva para otra cosa después).\\
\lstset{language=Matlab,frame=tb,framesep=5pt,basicstyle=\scriptsize}   
\begin{lstlisting}
function criptograma = CifrarTexto(llano,clave)
raiz = floor(sqrt(clave));
resto = clave-raiz^2;
N = length(llano);
for i=1:N,
    [raiz, resto, digito] = 
       DigitoDecimalRaizCuadrada(raiz,resto);
    cfr = llano(i)+digito;
    if cfr>90
        cfr = (cfr-91)+65;
    end
    criptograma(i) = char(cfr);
end
\end{lstlisting} 
 
\noindent

Fin de la introducción.\\


\begin{entradilla}
{\em Recomendación: {\color{introcolor}{para este tipo de texto}}
  utilizarlo para resaltar lo más importante o llamativo de la introducción}
\end{entradilla}



Linea de texto.

\sectiontext{white}{black}{II. Metodología}

Esta sección es recomendable para los artículos que contengan o hagan referencia a investigación, que incluyan experimentos.

\sectiontext{white}{black}{ Subsección de la Metodología}

Especifica brevemente elementos y cómo se realiza, en caso de que pueda replicarse.


\sectiontext{white}{black}{III. Nombre de la primera Sección}

Contenido del artículo. \\

\textsf {\textbf {Código original }}

(Esta sección solo es informativa, no se requiere que vaya en todos los artículos).\\ Código tomado de la revista Occam's Razor

{\href{https://groups.google.com/forum/#!forum/revista-occams-razor}{Grupo en Google}} \\
Disponible en:  
{\href{http://www.papermint-designs.com/roor/}{Revista online Occam's razor}}

Es mejor incluir imágenes, quita un poco el tedio y da más vida al texto para llamar la atención.

\begin{center}
\resizebox{5cm}{!}{{\epsfbox{images/nombre_articulo/glud.eps}}}
\mycaption{Gráfico del movimiento de los datos multimedia desde la generación del fuente hasta que lo escucha el usuario o oyente de la radio o canal de vídeo}
\end{center}

% Siguiente página
%%%%%%%%%%%%%%%%%%%%%%%%%%%%%%%%%%%%%%%%%%%%%%%%%%%%%%%%%%%%
\ebOpage{introcolor}{0.35}{Seccion dejar vacio}
%%%%%%%%%%%%%%%%%%%%%%%%%%%%%%%%%%%%%%%%%%%%%%%%%%%%%%%%%%%%

Para añadir un link:
\href{https://glud.org/}{Hipervínculo} \\
\href{https://blog.desdelinux.net/occamss-razor-publicacion-libre-de-divulgacion-cientifica/}{Acerca de Occam's razor} \\


Se repite la imagen anterior por flojera de cargar otra imagen, las imágenes que se agreguen en cada artículo deben ser diferentes del Kokopelli, dado que éste último estará presente en la portada de la revista. \\ 

Fin de la primera sección.



\sectiontext{white}{black}{IV. Sección con items}
Si se requiere enumeración o items:
\begin{itemize}
    \item Primero
    \item Segundo
\end{itemize}

\begin{enumerate}
    \item Primero
    \item Segundo
\end{enumerate}

Enhorabuena! 


\begin{entradilla}
{\em Este tipo de texto entradilla {\color{introcolor}{ (así se define en el código) }} 
  puede ir en cualquier parte del texto (fuera del resumen).}
\end{entradilla}

\sectiontext{white}{black}{V. Sección con tablas}\\
Si se requiere incluir tablas:
\medskip


\begin{tabular}{|l|l|l|l|}
    \hline
    & \multicolumn{3}{c|}{Europa} \\
    \cline{2-4}
    & Ciudad & Ró & S\'mbolo\\
    \hline \hline
    \multirow{3}{1cm}{España} & Madrid & Manzanares & Cibeles\\ \cline{2-4}
    & Sevilla & Guadalquivir & Giralda\\ \cline{2-4}
    & Zaragoza & Ebro & Pilar\\ \cline{1-4}
    Francia & París & Sena & Torre Eiffel\\ \cline{1-4}
    \multirow{2}{1cm}{Italia} & Roma & Tíber & San Pedro\\ \cline{2-4}
    & Milán & \multicolumn{1}{c|}{-} & Duomo\\ \cline{1-4}
\end{tabular}

\medskip

Otros tipos de tablas:
\href{https://latexlive.files.wordpress.com/2009/04/tablas.pdf}{Tablas en \LaTeX pdf} \\


\bibliographystyle{abbrv}
\begin{bibliografia}
\bibitem{kopka}
H.~Kopka and P.~W. Daly, \emph{A Guide to \LaTeX}, 3rd~ed.\hskip 1em plus
  0.5em minus 0.4em\relax Harlow, England: Addison-Wesley, 1999.
\bibitem{kopka1}
A.~Kopka and P.~W. Daly, \emph{A Guide to \LaTeX}, 3rd~ed.\hskip 1em plus
  0.5em minus 0.4em\relax Harlow, England: Addison-Wesley, 1999.
\bibitem{kopka2}
B.~Kopka and P.~W. Daly, \emph{A Guide to \LaTeX}, 3rd~ed.\hskip 1em plus
  0.5em minus 0.4em\relax Harlow, England: Addison-Wesley, 1999.
\bibitem{kopka3}
A.~Kopka and P.~W. Daly, \emph{A Guide to \LaTeX}, 3rd~ed.\hskip 1em plus
  0.5em minus 0.4em\relax Harlow, England: Addison-Wesley, 1999.
\bibitem{kopka4}
A.~Kopka and P.~W. Daly, \emph{A Guide to \LaTeX}, 3rd~ed.\hskip 1em plus
  0.5em minus 0.4em\relax Harlow, England: Addison-Wesley, 1999.
\bibitem{kopka5}
A.~Kopka and P.~W. Daly, \emph{A Guide to \LaTeX}, 3rd~ed.\hskip 1em plus
  0.5em minus 0.4em\relax Harlow, England: Addison-Wesley, 1999.
\bibitem{kopka6}
A.~Kopka and P.~W. Daly, \emph{A Guide to \LaTeX}, 3rd~ed.\hskip 1em plus
  0.5em minus 0.4em\relax Harlow, England: Addison-Wesley, 1999.
\end{bibliografia}


\begin{biografia}{images/nombre_articulo/glud.eps}{Nombre Completo del Autor} 
% añadir fotografía tamaño [2.5 cm x 3.3 cm ] 
%  |¬_¬|'  La foto es Opcional! , se añade reemplazando:
% \begin{biografia}{images/mi_articulo/autor.eps}{Nombre Completo del Autor}
Corta descripción del autor, que incluya estudios de pregrado o postgrado, pertenencia a grupos de investigación
\end{biografia}

\raggedcolumns
\pagebreak


\end{multicols}

\clearpage
\pagebreak

% Este fichero es parte del Número 5 de la Revista Occam's Razor
% Revista Occam's Razor Número 5
%
% (c)  2010 The Occam's Razor Team
%
% Esta obra está bajo una licencia Reconocimiento 2.5 España de Creative
% Commons. Para ver una copia de esta licencia, visite 
% http://creativecommons.org/licenses/by/2.5/es/
% o envie una carta a Creative Commons, 171 Second Street, Suite 300, 
% San Francisco, California 94105, USA.

% Este .tex contiene el contenido del articulo
% Seccion (Dejar en blanco)
%

% definiciones
\definecolor{filacolor}{HTML}{E33D00}
\definecolor{columnacolor}{HTML}{E65B05}


\rput(2.5,-2.3){\resizebox{!}{5.7cm}{{\includegraphics[]{images/Importancia/white.eps}}}}

% Cabecera
\begin{flushright}
\msection{introcolor}{black}{0.35}{Ingeniería de Software}

\mtitle{10cm}{Tendencias de la Ingeniería de Software en 2017.}

\msubtitle{8cm}{Trends in Software Engineering in 2017.}

{\sf Steven Sierra Forero, Ingeniero de Sistemas}\\


{\psset{linecolor=black,linestyle=dotted}\psline(-12,0)}
\end{flushright}

\vspace{2mm}
% -------------------------------------------------

\begin{multicols}{2}

\sectiontext{white}{black}{Resumen}

% Introducción
\intro{introcolor}{E}{ste paper expone de una manera generalizada las tendencias de la ingeniería de software en el 2017, aquí se expondrán los campos más relevantes en el desarrollo de software que están haciendo cambios significativos en los flujos de trabajo de las grandes compañías de desarrollo de software y por ende cambiando la forma en la que accedemos a la tecnología. \\
Metodologías como las contempladas por Agile, apoyando el desarrollo rápido de aplicaciones, despliegue de software automatizado en la nube con Kubernetes, y el uso intensivo de servicios en la nube como AWS en el ámbito de nubes privadas y empaquetamiento de software a través del uso de tecnologías de contención de software como Docker han sentado un precedente en la actualidad reevaluando constantemente la forma de hacer el delivery de productos de software y cambiando el pensamiento de quienes se involucran en el día a día de la ingeniería de software.
}

\sectiontext{white}{black}{Palabras clave: }Cloud computing, código abierto, ingeniería, software.

\vspace{5mm}

\sectiontext{white}{black}{I. Introducción}

En la actualidad la ingeniería de software ha tenido que enfrentar diversos cambios adaptándose a la necesidades de la industria, metodologías de planeación, construcción y evaluación del software han sido reevaluadas siguiendo el ritmo de la globalización. \\

A través de un recorrido general sobre el estado del arte de las tendencias en la ingeniería de software, se tratarán temas como el desarrollo de software en el campo del código abierto u open source, el desarrollo y uso de aplicaciones móviles, metodologías como las contempladas por Agile y finalmente computación en la nube, las cuales dan cuenta del avance de la ciencia y la tecnología en el siglo XXI los cuales cuales promueven cambios positivos para el desarrollo de la sociedad.

\vspace{5mm}

\sectiontext{white}{black}{II. OPEN SOURCE}

El software de código abierto ha tenido grandes avances en los últimos años, un referente muy conocido es RedHat, quien tomó el sistema operativo Linux y lo adaptó a las necesidades de un mercado corporativo, haciendo uso de su código de fuente abierta. El código abierto ha evolucionado desde los días en que Microsoft lo etiquetó como \textit{"un destructor de la propiedad intelectual"}, hoy en día una nueva generación de líderes de Microsoft espera que su plataforma Azure basada en Linux (con sus SDKs de código abierto) proporcione un crecimiento significativo a la compañía ya sus clientes. \\

Históricamente la tecnología open source, no ha sido respetada en el sentido de ser malinterpretada como software no estable o software sobre el cual no es viable hacer desarrollos corporativos de gran escala. Sin embargo en la actualidad existen organizaciones como Linux Foundation que tienen proyectos que no solo controlan Linux, sino también Cloud Foundry, el Proyecto Hyperledger (enfocado en la cadena de bloqueo empresarial), la Open Connectivity Foundation (que recientemente se fusionó con la antigua Alianza AllSeen y que se centra en la Internet de las Cosas ), La Fundación JS (centrada en JavaScript).\\

También la fundación Apache (Apache software foundation) mantiene la tradición de código abierto de software libre a través de su estricta licencia con una comunidad bastante activa. Actualmente, Apache tiene más de 180 proyectos, incluyendo nombres conocidos como CloudStack, Hadoop, Cassandra, Groovy, Kafka, Mesos, Maven y Tomcat.\\


Según la revista Forbes [3], en el evento anual que se lleva a cabo en los Estados Unidos, donde se exponen las últimas tendencias en el mercado del software de código abierto (open source) algunos de los proyectos más relevantes que se expusieron fueron:\\


% Siguiente página
%%%%%%%%%%%%%%%%%%%%%%%%%%%%%%%%%%%%%%%%%%%%%%%%%%%%%%%%%%%%
\ebOpage{introcolor}{0.35}{Ingeniería de Software}
%%%%%%%%%%%%%%%%%%%%%%%%%%%%%%%%%%%%%%%%%%%%%%%%%%%%%%%%%%%%

\begin{itemize}
\item Realidad aumentada y Realidad Virtual, esto se evidenció con Pokémon Go un juego para móviles que estremeció el mercado de las aplicaciones móviles con la integración de realidad aumentada y el uso de librerías de código fuente.

\item Inteligencia artificial la cual parecía el material de las películas no hace mucho tiempo. El aprendizaje no  supervisado hecho por la IA es un gran salto en el desarrollo de software. En este año se esperan sistemas que realmente puedan aprender y cambiar su comportamiento, abriendo el camino para dispositivos más inteligentes. Ejemplo de ello son todos los productos que ahora están en el mercado en liderado por los gigantes de la tecnología como Google, Amazon, Apple, Tesla e IBM quienes quieren potenciar el uso de inteligencia artificial en todos los productos de software que hacen parte del diario vivir de los consumidores de servicios digitales [1].
\end{itemize}





\sectiontext{white}{black}{III. APLICACIONES MÓVILES}

Los teléfonos inteligentes se han convertido en una parte inseparable de la vida cotidiana de muchas personas, empezando por la alarma en la mañana, el uso de WhatsApp, la música en la nube en aplicaciones como Spotify, videos en YouTube, correos electrónicos, con Gmail, Yahoo, Outlook, todas	estas	aplicaciones	que	se ejecutan sin problemas desde nuestros dispositivos móviles. Actualmente 2.1 billones de personas en todo el mundo posee un Smartphone, una persona en promedio usa su teléfono unas 264 veces al día, incluyendo textos y llamadas, es un número alto, pero corresponde a una realidad cotidiana en 2017, por lo tanto, el uso de Smartphone y aplicaciones móviles no se va a reducir en 2017 al contrario las tendencias de las aplicaciones para móviles pueden evidenciar nuevos horizontes en el desarrollo de software.\\

Tecnologías como las implementadas por Google en el proyecto AMP (Accelerated Mobile Pages) en donde Google ha anunciado que habrá un índice de búsqueda independiente para la web móvil abre un campo revolucionario que cambia completamente las tendencias de desarrollo de aplicaciones móviles, especialmente desde SEO y la perspectiva de aplicaciones web.
\\
Un ejemplo de las aplicaciones móviles desarrolladas en el ámbito corporativo puede ser Evernote a través del cual diferentes equipos en las compañías pueden colaborar en proyectos desde sus dispositivos móviles. Por otro lado, Facebook messenger y aplicaciones de listas de tareas son ejemplos de micro-aplicaciones que están soportando los procesos corporativos de las compañías en 2017.




\sectiontext{white}{black}{IV. METODOLOGÍAS ÁGILES}

En tanto que las metodologías Ágiles y DevOps no desaparecerán de ninguna manera, los líderes de software empresarial pasarán de tener visión de escalar en los procesos ágiles y hacer un esfuerzo  en la integración de profesionales DevOps ampliando la cadena de valor en el delivery de software de punta a punta. En 2017 hay una tendencia generalizada por la adopción de equipos de profesionales DevOps en procesos ágiles lo cual permite una maduración de las organizaciones en la entrega de software a un punto en el que están dispuestos a ver la agilidad de manera holística.\\

A medida que los líderes de software empresarial comienzan a enfocarse en su agilidad empresarial holística y en sus flujos de valor, buscan una solución definitiva en el reto de unificar la planificación estratégica , el desarrollo ágil y la implementación de DevOps.\\

Durante muchos años, las grandes organizaciones empresariales se han centrado en acelerar la entrega escalando ágil. Muchas de estas organizaciones también se han centrado en mejorar la planificación estratégica inicial con prácticas ágiles. En los últimos años estas organizaciones han comenzado a centrarse en los beneficios que DevOps puede tener para la entrega de productos de software en un menor tiempo, lamentablemente, cada una de estas iniciativas se realiza típicamente de forma aislada dando como resultado un esfuerzo fragmentado que no llega a toda la organización.\\

La mayoría de los flujos de valor de las organizaciones empresariales se dividen en tres fases principales [3]: estrategia, desarrollo y entrega. En un flujo de valor fragmentado, cada fase evoluciona en sus propios equipos, procesos, herramientas e información que no están totalmente integrados con el resto del flujo de valor. Esta evolución ha creado un entorno en el que, aunque la adopción ágil y de DevOps es alta, los resultados no satisfacen las expectativas de liderazgo.\\

Una buena cantidad de organizaciones empresariales reconocidas se han esforzado en los últimos años en el mejoramiento del proceso ágil optimizando los flujos de valor de toda la empresa [2].\\

Esta es una transformación importante que no es fácil de lograr. Afortunadamente, la lección aprendida por grandes organizaciones es un punto de partida para otras organizaciones que están empezando en el mercado de la ingeniería de software.\\

% Siguiente página
%%%%%%%%%%%%%%%%%%%%%%%%%%%%%%%%%%%%%%%%%%%%%%%%%%%%%%%%%%%%
\ebOpage{introcolor}{0.35}{Ingeniería de Software}
%%%%%%%%%%%%%%%%%%%%%%%%%%%%%%%%%%%%%%%%%%%%%%%%%%%%%%%%%%%%


\sectiontext{white}{black}{V. CLOUD COMPUTING}

La computación en nube ha reducido drásticamente el costo del despliegue de software. En lugar de necesitar equipos de cómputo y hardware especializado, espacio en el centro de datos,  nómina de administradores de sistemas, el software se puede contratar de los proveedores de nube con unos pocos clics y un costo relativamente bajo por hora. El éxito ya no es un problema, ya sea: construya bien su software y los recursos adicionales de la nube le ayudarán a escalar hacia arriba y hacia abajo según sea necesario. Las pequeñas empresas e incluso los ingenieros de software individuales pueden proporcionar el mismo nivel de servicio en la nube que las grandes empresas.
El porcentaje de empresas en la actualidad que tienen una estrategia para utilizar despliegue [6] de software en la nubes creció hasta el 85\% (vs 82\% en 2016), con un 58\% de planificación híbrida (vs. 55 por ciento en 2016). También hubo un aumento en el número de empresas que planean múltiples nubes públicas (de un 16\% a un 20\%) y una disminución concurrente en las que planean múltiples nubes privadas (de 11\% a 7\%). En 2017, los desafíos para el despliegue de software en la nube se redujeron en todos los ámbitos, con la excepción de la gobernanza y el control, que se mantuvo estable [3].\\

Como parte de la adopción de los procesos de DevOps para el despliegue de software en la nube, las empresas han optado por implementar nuevas herramientas que les permitan estandarizar y automatizar la implementación y configuración de servidores y aplicaciones. Estas herramientas incluyen herramientas de gestión de configuración (como Chef, Puppet y Ansible),  más recientemente, tecnologías de contenedores, como Docker y orquestación de contenedores y herramientas de programación como Kubernetes, Swarm y Mesosphere.\\

La encuesta sobre el estado de la nube de 2017 realizado por la Cloud Native Computing foundation revela que aunque AWS continúa liderando la adopción de la nube pública (57\% de los encuestados actualmente ejecutan aplicaciones en AWS), este número se ha mantenido plano desde 2016 y 2015. Es importante señalar que mientras que el porcentaje de las empresas que ejecutan al menos una aplicación en AWS es plana, el número de aplicaciones y máquinas virtuales que están ejecutando está aumentando, con lo que el aumento de los ingresos de AWS. El aumento sustancial en el uso de contenedores ahora hace Docker la herramienta de DevOps superior entre los incluidos en nuestra encuesta. La adopción general de Docker subió al 35\%, tomando la ventaja sobre Chef y Puppet en 28\%.\\

La computación en nube ha ayudado a muchas empresas a transformarse en los últimos cinco años, pero los expertos están de acuerdo en que el mercado está entrado en una nueva etapa, tanto para cloud público como para cloud computing construido o alojado en centros de datos corporativos. El mercado de la nube se ve creciendo a un buen ritmo en en 2017 y lo seguirá haciendo medida que las empresas busquen obtener eficiencias a medida que escalan sus recursos computacionales para servir mejor a los clientes.\\


Según Dave Bartoletti, "Las empresas con grandes presupuestos, centros de datos y aplicaciones complejas ahora están viendo a cloud como un lugar viable para ejecutar aplicaciones empresariales básicas". Esto se evidencia en las estadísticas entregadas por Amazon donde se evidencia que con la primera ola de cloud computing creada como Amazon Web Services, en 2006 ha dado utilidades por más de 11.000 millones de dólares.\\


\bibliographystyle{abbrv}
\begin{bibliografia}
\bibitem{ref1}
\href{http://insights.dice.com/trends/}{insights.dice.com}
\bibitem{ref2}
\href{https://www.oreilly.com/ideas/5-software-development-trends-shaping- enterprise}{www.oreilly.com}
\bibitem{ref3}
\href{https://forbes.com}{forbes.com}
\bibitem{ref4}
Barry W. Boehm. Get ready for the agile methods, with care. IEEE Computer, 35(1):64–69, Jan 2002.
\bibitem{ref5}
Alan C. Wills. Agile components. Technical report, Trireme International	Ltd,	Feb	2002. \href{http://www.ltt.de/otland/experts/a.c.wills.shtml}{www.ltt.de}
\bibitem{ref6}
Mihaly Csikszentmihalyi. Flow: The Psychology of Optimal Experience. Harper Perennial Modern Classics. 2008.
\bibitem{ref7}    
Jeannette M. Wing. Computational thinking. Communications of the ACM, March 2006. Page 33- 35.
\end{bibliografia}


\begin{biografia}{images/Importancia/white.eps}{Steven Sierra Forero} Estudiante de Ingeniería de Sistemas, Universidad Distrital Francisco José de Caldas.\\ssierraf@correo.udistrital.edu.co
\end{biografia}

\raggedcolumns
\pagebreak


\end{multicols}

\clearpage
\pagebreak

% Este fichero es parte del Número 5 de la Revista Occam's Razor
% Revista Occam's Razor Número 5
% (c)  2010 The Occam's Razor Team
% Esta obra está bajo una licencia Reconocimiento 2.5 España de Creative
% Commons. Para ver una copia de esta licencia, visite 
% http://creativecommons.org/licenses/by/2.5/es/
% o envie una carta a Creative Commons, 171 Second Street, Suite 300, 
% San Francisco, California 94105, USA.

% definiciones
\definecolor{filacolor}{HTML}{E33D00}
\definecolor{columnacolor}{HTML}{E65B05}

% Incluye imagen del artículo (Debe ser diferente del Kokopelli! (°_°) Sin embargo, no es obligatoria esta imagen)

\rput(2.5,-2.3){\resizebox{!}{5.7cm}{{\includegraphics[]{images/radioglud/entrada.eps}}}}


%Imagen de el comienzo de el articulo, coordenadas desde %la parte superior izquierda del margen de la pagina
% Las imagenes quedan mejor en formato .eps, agradecería si alguien logra colocar imagenes en dicho formato y comparte la corrección.
% Cabecera
\begin{flushright}
\msection{introcolor}{black}{0.35}{Tecnologías Libres}

\mtitle{10cm}{Emisora en línea con Tecnologías Libres}

\msubtitle{8cm}{Online Broadcasting Station with Free and Open Source Tecnologies}

{\sf Jorge Useche, Ingeniero Electrónico}\\
{\sf Alejandro Cortázar, Bachiller Académico}

{\psset{linecolor=black,linestyle=dotted}\psline(-12,0)}
\end{flushright}

\vspace{2mm}
\begin{multicols}{2}

\sectiontext{white}{black}{Resumen}

% Introducción
\intro{introcolor}{E}{n nuestra sociedad la radio hace aún parte importante de nuestras vidas; tal vez escuchada más por personas que han venido envejeciendo con estos medios y desplazada por el Internet una nueva forma más eficiente de obtener información de todo tipo sin necesidad de publicidad ni cosas innecesarias porque siempre hay muchas opciones de donde escoger. 

Es necesario recuperar la calidez humana que nos brinda escuchar la información de otra persona, imaginar de cierta forma y creer en un mundo que es generado solo por una voz; aprovechar el Internet para que la gente de distintos rincones accedan a la información y a los pensamientos que queremos compartir.
El espíritu de compartir se ve truncado en veces por no acceder a las tecnologías por ser muy caras y ahí es donde entra el software libre a permitirnos compartir lo que queramos, y para la intensión anterior existe un conjunto de herramientas que permitirán que tengamos nuestra propia emisora en Internet o netradio, muy fácil de manejar y gratuita.
}

\sectiontext{white}{black}{Palabras clave: } Radio, Online, Comunitaria.

\sectiontext{white}{black}{I. Introducción}

% Cuerpo del artículo
Siempre es mejor confiar en la estabilidad de un sistema GNU/Linux que en la de otros sistemas operativos; para este fin se recomienda ampliamente utilizar una distribución como Debian GNU/Linux o algo aún más libre como Trisquel GNU/Linux ya que nos ofrece paquetes de excelente calidad y muy estables. Por esta razón aparte de la cuestión moral de que sea libre, únicamente se utilizará tecnologías y contenidos libre que cumpla las 4 libertades del software según Richard Stallman adaptadas al tipo de contenido. La creación del sistema para la emisora en Internet se abordará con el siguiente método:\\

\begin{enumerate}
\item ¿Por qué una radio en línea?
\item Conceptos generales sobre servicios o demonios en GNU/Linux.
\item Terminología usada en la transmisión de medios.
\item Piezas de software usadas.
\item Alternativas para emisión de audio y ventajas y desventajas de éstas.
\item ¿Qué música puedo poner?, licencias asociadas.
\item Ilustración general del método usado en Radio GLUD.
\item Instrucciones para la primera emisión en una red local.
\end{enumerate}

Con esto claro se entienden ya los alcances de este escrito y se quiere dejar en claro, que para entender lo que se hace para crear la emisora en línea es necesario que se expliquen los conceptos relacionados y la razón de todo lo que hallamos tenido que decidir, para que entendiendo esto se pueda automatizar procesos o mejorar la forma en que se hicieron, de lo contrario el artículo que convertiría en un tutorial más.

\begin{entradilla} %codigo para una entradilla
{\em {\color{introcolor}{Radio GLUD}} \\ Libertad para tus oídos.}
\end{entradilla}

Antes de desarrollar los contenidos, me parece necesario decir porque se cree tener el criterio necesario para escribir este artículo. Desde hace un tiempo existe la emisora en línea del GLUD, esta emisora se llama ``Radio GLUD'', y desde sus comienzos ha dado valioso conocimiento a todo el grupo con el cual se han solucionado algunos problemas con soluciones definitivas o parciales en que a pesar de todo el trabajo invertido hasta ahora, todavía falta un camino largo por recorrer para acercarla lo más posible a lo que puedan imaginar nuestras mentes.
La invitación cordial va a que visiten el sitio web de este proyecto \url{http://glud.udistrital.edu.co/radioglud}

% Siguiente página
\ebOpage{introcolor}{0.35}{Tecnologías libres}

\sectiontext{white}{black}{II. ¿Por qué una radio en línea?}

Es una pregunta que se hizo a su debido momento, por lo general se presentan muchas cosas a favor y otras en contra, pero las razones más valiosas que se encontraron fueron positivas y más enfocadas en {\em impacto social} que en cuestiones técnicas. 



La madre de todas es el poder social que se presenta cuando las personas comparten algo, pero ¿cuál era el medio correcto para compartir?; la radio es nuestro ancestro más cercano (para comunicarnos), todavía vivo, más interesante y nos ofrece las siguientes ventajas derivadas del poder social.

El primero de ellos es {\em el sentido de comunidad}, que se crea cuando muchas personas trabajan por una misma causa y es conocido que es mucho mas fácil hacer un programa radial que uno televisivo; el segundo es {\em compartir ideas y sentimientos}, que se ha visto mucho en otras radio, pero ¿realmente son las que nos interesan en la comunidad académica?; la tercera es {\em las experiencias que se adquieren}, no puede ser solo libros, pueden haber otros espacios de esparcimiento muy provechosos; el cuarto es {\em compartir gustos musicales}, que hace que estemos más unidos, o ¿quién no ha cantado una canción que a todo mundo le gusta con un grupo de amigos?; y la última, {\em la creatividad como reflejo de sociedad}, tenemos derecho a expresarnos y nos volvemos más creativos (A mi parecer) solo mostrando nuestra voz que nuestros rostros.

\sectiontext{white}{black}{III. Conceptos generales sobre servicios o demonios en GNU/Linux}

Primero abarcaré la forma en que el sistema operativo gestiona programas como icecast (ver sección Piezas de software usadas). Dicho programa se gestiona como un servicio (daemon en sistemas unix) donde la principal caracteristica es que es un programa que es persistente con una ejecución continua, se ejecuta en segundo plano y no posee una interfaz de usuario directa o shell.

Disk And Execution MONitor es lo que significa DAEMON; ``según una investigación realizada por Richard Steinberg, la palabra fue utilizada en 1963 por primera vez, en el área de la informática, para denominar a un proceso que realizaba backups en unas cintas, este proceso se utilizó en el proyecto MAC del MIT y en una computadora IBM 7094,1 dicho proyecto estaba liderado por Fernando J. Corbató, quien afirma que se basó en el demonio de James Maxwell''\cite {ref1}

Hay dos formas de iniciar un servicio si éste no se encuentra activo, la primera es ejecutando el programa relacionado con los parámetros para iniciarlo como daemon que indique el manual y la segunda es por medio de un mediador como service (programa en shell) que facilita la tarea ya que no hay que recordar la ruta ni el comando de ejecución del servicio.

\sectiontext{white}{black}{IV. Terminología usada en la transmisión de medios}

La transmisión de medios específicamente de audio, representa la acción esencial del proceso de crear una emisora, hay que tener claro que el proceso no es más que el streaming de audio que posee una interfaz web para poder escucharlo cómodamente.

\begin{center}
\resizebox{8cm}{!}{{\epsfbox{images/radioglud/mov-dat-mult.eps}}}
\mycaption{Gráfico del movimiento de los datos multimedia desde la generación del fuente hasta que lo escucha el usuario o oyente de la radio o canal de vídeo}
\end{center}

El streaming (que se refiere a corriente continua) es la distribución de contenido multimedia a través de una red de computadores con la ventaja de que por medio de un búfer de datos se agrega la posibilidad de disfrutar el contenido mientras este descarga; el componente búfer de datos en este caso es un espacio de memoria que almacena datos para evitar que la reproducción del medio se quede sin información durante la transferencia.


Latencia (relay - retardo del reloj) es el tiempo o retardos temporales en que se realiza un proceso o procesos (o en su defecto ninguno) debido por lo general al tiempo de propagación de los circuitos o incluso al tiempo en que se estabiliza que es debido a la densidad del material y la temperatura principalmente.

\begin{center}
\resizebox{8cm}{!}{{\epsfbox{images/radioglud/latencia.eps}}}
\mycaption{Gráfico de un sistema intermedio (sistema 2) donde se muestra que existe un tiempo desde que los datos entran hasta que salen del sistema y los causales posibles de ese Delta}
\end{center}


% Siguiente página
%%%%%%%%%%%%%%%%%%%%%%%%%%%%%%%%%%%%%%%%%%%%%%%%%%%%%%%%%%%%
\ebOpage{introcolor}{0.35}{Tecnologías libres}
%%%%%%%%%%%%%%%%%%%%%%%%%%%%%%%%%%%%%%%%%%%%%%%%%%%%%%%%%%%%



\begin{center}
\resizebox{8cm}{!}{{\epsfbox{images/radioglud/propagacion.eps}}}
\mycaption{Gráfico de los tiempos de propagación de un sistema, comenzando con el tiempo de propagación de bajo a alto (tplh) y terminando su respuesta con el tiempo de propagación de alto a bajo (tphl)}
\end{center}

El tiempo real (Real Time - RT), se presenta esta referencia para todos los sistemas en donde hay que considerar que ese tiempo real no tiene que ser necesariamente sinónimo de rapidez o inmediatez; como concepto, dicho tiempo es la latencia suficiente para que el sistema pueda resolver el problema para el cual está dedicado. Para el caso del audio en informática existe un tiempo de respuesta crítico con un porcentaje que no exceda el 3\% (aunque por lo general es así con todos los sistemas RT) que para llevarse a cabo necesita algunos componentes esenciales como un sistema operativo (OS) que soporte el tiempo real (para el caso de un sitema GNU/Linux solo sería el kernel) y un software diseńado para esto (el mismo OS RT ofrece un marco para el desarrollo de dichas aplicaciones).

Hay que distinguir entre 2 tipos de sistemas RT, los pasivos y los activos donde la diferencia es el dańo o fallo que puede causar un retraso en la seńal; por ejemplo tenemos que un sistema de audio o vídeo en directo sería un sistema RT pasivo debido a que un retraso en la seńal simplemente será un retraso en la visibilidad del vídeo o escucha del audio tal vez con desmejora del multimedia pero luego seguirá funcionando normal como si nada hubiera pasado y un ejemplo de un sistema RT activo son la mayoría de sistemas médicos (como el marcador de pasos artificiales o marcapasos) donde un retraso de la seńal puede causar el fallo del entorno donde el sistema interfiere causando una desmejora apreciable del paciente o incluso la muerte.

El servidor de sonido para streaming es un software con la función de gestionar los flujos de datos 
reproduciendo las muestras. Cualquier programa reproductor, grabador, gestor de audio, necesita soporte para trabajar con un determinado servidor de audio por ejemplo Jack Audio Connection Kit (en adelante simplemente JACK) es uno de los más conocidos para flujos en tiempo real.

El servidor de streaming de medios es un sistema de distribución de multimedia por streaming que utiliza un tipo de formato y un protocolo determinados para realizar el proceso de distribución de contenidos multimedia, para el caso de icecast el formato puede ser ogg, webm, aac o mp3 y el protocolo http o https.

El flujo multimedia se refiere al audio o al vídeo que es enviado desde un punto de la red hasta otro. Este flujo debe ser continuo aunque no necesariamente estable (fluctuante), pero debe tener un mínimo de velocidad de transferencia de un punto a otro para que el streaming sea efectivo, hay que introducir entonces algunos conceptos relacionados a las propiedades intrínsecas del archivo multimedia que se quiere enviar por streaming para estimar esa velocidad mínima que debe tener la red. Hay que tener en cuenta que cualquier archivo multimedia (audio y/o vídeo) posee un bitrate, que es la cantidad de bits que posee el flujo en unidad de tiempo, y que obedece a la composición del archivo, se mide por lo general en kbps (kilo bit por segundo).

También es importante conocer el término multi-bitrate que surgió de la necesidad de enviar una fuente multimedia a diferentes calidades, por ejemplo una calidad HD, normal y móvil; para cada una de estas es necesario un bitrate distinto para que el usuario pueda consumirlo de acuerdo a sus necesidades o limitación de ancho de banda de su conexión.

Ya que introducimos los bitrates y las velocidades de conexión de la red, es necesario entrar un poco a hacer el análisis de requerimientos de ancho de banda. Cuando queremos enviar un flujo multimedia desde nuestro computador hasta un servidor de streaming es pertinente revisar la velocidad de conexión (de subida) a este servidor para enviar nuestro flujo o flujos; para llevar a cabo este análisis es necesario realizar  un test de velocidad o Speed Test para saber nuestra velocidad real de conexión; hay que tener en cuenta que pueden presentarse fluctuaciones inesperadas en esa velocidad máxima por lo tanto es recomendable enviar flujos con bitrates que sumados estén dentro del rango de un 50\% a un 70\% de la capacidad total, una vez tengamos el dato podemos saber cuantos flujos y de que calidades podemos enviar a el servidor; por ejemplo si nuestra conexión es de 512 kbps se deben hacer los cálculos sobre la base de 358,4 kbps que es el 70\% de la conexión, con esto se puede enviar audio con bitrates de 256 kbps y 64 kbps sin mucho problema.

Source Stream (fuente) es entendido como la fuente de audio (o vídeo) de donde se obtienen los datos para despues distribuirlos en la red en forma de streaming, este fuente se genera desde un cliente como ices o idjc de los que se hablará más adelante.

Mount Point (Punto de Montaje) es una ruta por la cual se puede acceder a determinado streaming, debe ser configurada en el servidor de streaming, un ejemplo de un punto de montaje es ``/radioglud.ogg'' y para acceder a él desde un navegador web o un reproductor de audio la dirección se escribe de la forma http://ip:puerto/radioglud.ogg, no hay que preocuparse si no se ha entendido del todo ya que al hacer nuestro proceso de streaming en localhost el concepto llega sin mucho esfuerzo.


% Siguiente página
%%%%%%%%%%%%%%%%%%%%%%%%%%%%%%%%%%%%%%%%%%%%%%%%%%%%%%%%%%%%
\ebOpage{introcolor}{0.35}{Tecnologías libres}
%%%%%%%%%%%%%%%%%%%%%%%%%%%%%%%%%%%%%%%%%%%%%%%%%%%%%%%%%%%%




También vale la pena saber algunos conceptos de envío de información conocidas como {\em formas de transmisión de información} en una red, en lo que se consideran 3 como los principales conceptos, hay que decir que existen más pero éstos son los más relevantes y usados, el primero es {\em unicast} que consiste en que se envía la información uno a uno, donde obtiene la información el que lo necesita a través de requerimientos; el siguiente es {\em multicast} que necesita un rango de IPs ya reglamentado (conocidas como tipo D) para enviar simultáneamente los datos a todos los destinos escogidos a los que se les conoce como grupo multidifusión, no necesita enviar la información a cada destino por petición, sino que realiza una única transmisión de los fuentes a cada uno de los destinos y por eso no necesita una red con gran velocidad para realizar las transferencias; y el último de los tres es la forma {\em broadcast} que es parecido al anterior con la diferencia que se envía los datos a todos los puntos de la red o subred simultáneamente.

Icecast es unicast y hasta donde conozco no tiene soporte multicast, uno de los indicios principales para esta conclusión es el protocolo HTTP (Protocolo de Transferencia de HiperTexto) que es TCP (Protocolo de Control de Transmisión) con el que icecast distribuye el streaming, además de que supongo que como he leído desde hace un tiempo a pesar de que la multidifusión tiene muchas ventajas sobre el bandwidth (ancho de banda de la conexión de Internet) su tráfico es realmente complejo además de que su implementación a gran escala produce puntos de fallos que pueden ser utilizados para ataques de denegación de servicios. Como dato adicional solo los protocolos diseńados específicamente para trabajar con multicast son efectivos para su implementación y la mayoría de las aplicaciones conocidas lo hacen por UDP (Protocolo de Datagrama de Usuario).

\begin{entradilla} %codigo para una entradilla
{\em {\color{introcolor}{Poder Social}} se presenta cuando las personas comparten.}
\end{entradilla}

\sectiontext{white}{black}{V. Piezas de software usadas}

A continuación se listan las aplicaciones que se han usado en Radio GLUD a través de su desarrollo, más adelante se dirá cuales de estas se usan actualmente.

\begin{itemize}
\item JACK: Es un servidor de sonido (daemon) de baja latencia, GNU GPL.
\item Icecast: Es un proyecto y a la vez un programa servidor de streaming de medios, creado por la Xiph.Org Foundation, su propósito es crear un flujo de medios para una estación de radio, jukebox privado (rockola) y otros, a partir de un source client, GNU GPL v2.
\item Ices: Es un programa generador de fuentes (source client) para un servidor de streaming, su propósito es generar un flujo de audio (stream) para el servidor de streaming del proyecto ICECAST, GNU GPL.
\item IDJC: Cliente generador de fuentes para un servidor de streaming, es mucho más versátil que ices ya que permite mezclar distintos flujos de audio, provenientes de cualquier dispositivo o programa (incluso de si mismo) que pueden estar gestionados principalmente por el servidor de audio local. Su función es principalmente servir de estudio de audio casero, GNU GPL v2.
\item QjackCtl: Es una aplicación QT, que nos da la posibilidad de controlar algunos parámetros de JACK (audio server daemon) de una forma simple y gráfica, GNU GPL v2.
\item JACK Rack: Es un programa con un conjunto de efectos para JACK a traves de complementos (plugins) su procesamiento del audio es a baja latencia, GNU GPL.
\item JAMin: Es una interfaz muy poderosa de compresión y ecualización de audio, GNU GPL.
\item Liquidsoap:  Es un lenguaje de programación de audio para generar source streams, al ser un interprete con un lenguaje script dedicado la versatilidad y la adaptabilidad lo hace lo suficientemente poderoso para acomodarse a cualquier emisora, por complejo que se quiera realizar el(los) source stream(s), GNU GPL.
\item jme (jmediaelement): Es una aplicación web, reproductor de audio y vídeo a través de elementos de HTML5.
\item Otras aplicaciones de audio muy útiles en  \url{http://jackaudio.org/applications}.
\end{itemize}

\begin{entradilla} %codigo para una entradilla
{\em {\color{introcolor}{Jack Audio Connection Kit}} multiplataforma y lo mejor que el software libre ha podido hacer por el mundo de la música.}
\end{entradilla}


% Siguiente página
%%%%%%%%%%%%%%%%%%%%%%%%%%%%%%%%%%%%%%%%%%%%%%%%%%%%%%%%%%%%
\ebOpage{introcolor}{0.35}{Tecnologías libres}
%%%%%%%%%%%%%%%%%%%%%%%%%%%%%%%%%%%%%%%%%%%%%%%%%%%%%%%%%%%%

\sectiontext{white}{black}{Alternativas para emisión de audio y ventajas y desventajas de estas}

Son muchas las aplicaciones que se utilizan en el entorno comercial para la emisión de medios, aunque contamos con pocas opciones de software libre, todas son de gran calidad y muy estables. Se listan las alternativas privativas que conozco como para dar ciertas observaciones al respecto, las que no conozco debido a que nunca las he usado, ni escuchado, ni desearía hacerlo aparecen con un ``--''.




\begin{center}
\begin{tablehere}
\begin{tabular}{|>{\columncolor{columnacolor}} c |>{\columncolor{columnacolor}} c |}
\hline
\multicolumn{2}{|>{\columncolor{filacolor}}c|}{Semi-equivalentes entre software libre y privativo}\\
\hline
\rowcolor{filacolor}Alternativa libre & Alternativa privativa \\ \hline 
Icecast & SHOUTcast\\ \hline
Icecast & oddcast o edcast\\ \hline
Ices, Darkice & Winamp, VLC\\ \hline
IDJC & Winamp, Zararadio\\ \hline
QjackCtl & --\\ \hline
JACK Rack & --\\ \hline
JAMin & Winamp\\ \hline
liquidsoap & --\\ \hline
jme (Reproductor HTML5)& Reproductor en Flash\\ \hline
WebcamStudio & WebcamMax\\ \hline
\end{tabular}
\caption{Equivalentes lo más parecido posible a las herramientas libres}
\end{tablehere}
\end{center}


Hace algún tiempo vi una conferencia en el CPCO4 (Campus Party Colombia), la verdad no recuerdo el nombre de la persona y no la encontré por la web por eso no la cito, en esta se enseńó como hacer una emisora online con software privativo. Algunas de las quejas del expositor en la conferencia eran, ``shoutcast solo funciona con winamp''; ``winamp solo trabaja con la tarjeta de audio escogida como principal''; ``la configuración de archivos de texto con notepad es feo (eso puede ser debido a que pone el texto todo de un solo color y posiblemente no se aprecia la indentación como debe ser)'' y dijo también ``estoy más perdido que el hijo de limber'' y por las constantes fallas en el proceso decía ``eso le pasa al chavo, al coyote y a mí''. Mis observaciones fueron que el equipo se colgaba fácilmente, la razón es que posiblemente el software era muy pesado, empezando por el sistema operativo.\\

Una de las cosas que si se hacían bien con el software privativo era cuando usaba herramientas como oddcast y zararadio aunque en sus versiones gratis poseen algunas limitaciones en cuanto a funcionalidad. Otra aplicación que usaba era WebcamMax que sirve para generar un fuente de vídeo con imágenes u otros vídeos para posteriormente unirla con el audio y enviarla a un canal en {\em ustream}, buscando un rato en la web encontré la alternativa libre {\em WebcamStudio} aunque por estar escrita para Java 1.6 ha resultado hasta el momento imposible hacerlo funcionar con nuevas versiones de Java, solo queda esperar a que la versión alpha siga desarrollándose y llegue a ser estable.\\

Queda rescatar que tu libertad está sobre todas las cosas, así que aunque algunas veces el software privativo sea más funcional no hay que usarlo porque luego te esclavizaras a los desarrolladores de dicha herramienta.


\sectiontext{white}{black}{VI. ¿Qué música puedo poner?, licencias asociadas}


Cuando hacemos nuestra propia emisora en línea es necesario conocer algunas leyes colombianas e internacionales que nos impiden poner música de cualquier tipo sin considerar que existen los derechos de autor (patrimoniales y conexos) y que hay que posiblemente pagar por ellos a un ente recaudador de derechos como Sayco y Acinpro (al parecer va a ser cambiada por otro ente debido a la corrupción del mismo), de lo contrario se está cometiendo un delito.\\
Los derechos de autor tienen distintas clases, entre los que interesan para el desarrollo de la pregunta están:
\begin{itemize}
\item Los derechos Morales: son los derechos ligados al autor y perduran por la eternidad, estos derechos son irrenunciables e imprescriptibles (ni ceder, vender o transferir).
\item Los derechos Patrimoniales: son los derechos que definen la explotación de la obra (retribuciones por uso, reproducción, difusión) hasta cierto límite después de la muerte del autor, después de este plazo se convierte en lo que se llama dominio público.
\item Los derechos Conexos: son los derechos que cobijan a personas distintas al autor (artistas, interpretes, traductores, editores, productores y otros)
\end{itemize}

\begin{center}
\begin{tablehere}
\begin{tabular}{|>{\columncolor{columnacolor}} c 
		|>{\columncolor{columnacolor}} c 
		|>{\columncolor{columnacolor}} c 
		|>{\columncolor{columnacolor}} c |}\hline
\multicolumn{4}{|>{\columncolor{filacolor}}c|}{Comparativa de tarifas de Sayco y Acinpro}\\ \hline
\rowcolor{filacolor}Lugar & Categoría & Estrato & Precio \\ \hline
Piojó 		& 6 & 1 & \$278.533   \\ \hline
Natagaima 	& 4 & 2 & \$371.472   \\ \hline
Ibagué 		& 2 & 2 & \$566.700   \\ \hline
Bogotá 		& E & 3 & \$1.133.400 \\ \hline
\end{tabular}
\caption{Datos tomados de \url{http://www.saycoacinpro.org.co/test_tarifas.php} con el código 102 (Establecimientos de cualquier naturaleza que para fines de su objeto social vendan productos para llevar y consumir o usar en sitios diferentes al del establecimiento comercial.Ejemplo: Almacenes de cadena, Supermercados, Hipermercados, Bancos, Estación.) y Capacidad/Aforo de 50 personas}
\end{tablehere}
\end{center}

% Siguiente página
%%%%%%%%%%%%%%%%%%%%%%%%%%%%%%%%%%%%%%%%%%%%%%%%%%%%%%%%%%%%
\ebOpage{introcolor}{0.35}{Tecnologías libres}
%%%%%%%%%%%%%%%%%%%%%%%%%%%%%%%%%%%%%%%%%%%%%%%%%%%%%%%%%%%%


Tal vez para una emisora libre de publicidad o una emisora estudiantil, sería muy difícil invertir recursos para la música y tener que suspender el uso de piezas musicales cuando es bien conocido que la música (ver arriba los argumentos) hace parte importante del proceso de comunidad que es lo que se busca en un grupo como GLUD.\\

Existen algunos tipos de productos musicales debidos al cumplimiento de los derechos anteriormente nombrados y que se podrán utilizar en nuestra emisora teniendo en cuenta que uso es el que se le va a dar
\begin{itemize}
\item Música libre: Es la música que cumple los mismos principios del software libre.
\item Música con algunos derechos reservados: Es la música que no permiten las libertades de obra derivada y/o uso comercial.
\end{itemize}

Dentro de la música que se puede usar está la de Dominio Público, mucha de la mejor música del planeta se encuentra bajo esta licencia, entre ellas obras de arte como la música clásica; también encontramos piezas de audio que el autor ha decidido registrar con licencias como Creative Commons, Coloriuris o ArtLibre.\\

Las licencias más usadas actualmente para contenidos libres o semilibres son las Creative Commons que son 6 licencias que surgen de la conjugación de las características o funciones que son, Atribución (by), No comercial (nc), Sin derivar (nd) y Compartir igual (sa). Las licencias libres\cite {ref2} son CC-by y CC-by-sa, las otras 4 que son CC-by-nd, CC-by-nc, CC-by-nd-nc y CC-by-nc-sa son las que se consideran semi-libres o con algunos derechos reservados. Pueden buscar su licencia ideal en la dirección \url{http://creativecommons.org/choose/} y posteriormente poner los meta-datos en tu sitio web.\\

Estas licencias CC (Creative Commons) según cuenta la historia surgió debido a la aparición del Internet (hay que darse cuenta la importancia de esta en lo que que consideramos que es la libertad por lo menos en la tecnología) y el espíritu de compartir que hicieron que una persona Eric Eldred fuera subiendo a la red obras que quedaban en Dominio Público cada año, en la escena apareció Mary Bono viuda de Sonny Bono un actor estadounidense que quería que los derechos de autor durarán para siempre, y luego de su muerte, Mary Bono tomó el cargo que él tenía en la cámara de representantes y comenzó a promover lo que ella consideraba la última voluntad de su marido.\\

Todo este lío resulta ser anticonstitucional pero igual se aumentó los ańos que se protegía el derecho patrimonial del autor, situación que causó descontento (para Eric Eldred y sus seguidores) ya que para ese ańo y otros venideros no habría liberación de obras; en esté dilema aparece Lawrence Lessig con un argumento contundente ``en la constitución de USA decía que se garantizaba el derecho monopólico por un tiempo limitado para asegurar y promover el progreso'', hubieron discusiones hasta que lo denominado ``The Mickey Mouse Protection Act''\cite {ref3} impuso con su poderío una influencia lo bastante grande para que el tiempo del derecho patrimonial se extendiera sin regulación aparente cada vez que llegaba la hora de que Mickey Mouse pasara a dominio público\cite {ref4} en 1976 decidieron parar estos cambios casi anuales. \\

Mickey Mouse ahora se esconde bajo una marca registrada que impide usarlo por ser prácticamente el logo de la compańía The Walt Disney Company y en el transcurso de todo esto Lawrence Lessig creó lo que actualmente es la organización sin fines de lucro Creative Commons como oposición directa al copyright.

\begin{entradilla} %codigo para una entradilla
{\em {\color{introcolor}{Internet}} y el espíritu de compartir!!!.}
\end{entradilla}
\end{multicols}

\sectiontext{white}{black}{Ilustración general del método usado en Radio GLUD}



\begin{center}
\myfig{0}{images/radioglud/radioglud.eps}{0.8}
\mycaption{Gráfico del método usado en Radio GLUD, consiste en 2 computadores donde el client streaming que genera el source stream es Mafalda GLUD y el servidor de streaming es glud.udistrital.edu.co}
\end{center}

\begin{multicols}{2}



Como se puede ver, el método (ha de entenderse método como el esquema general o a grandes rasgos de la gestión, configuración y flujo del audio) implementado hasta ahora no incluye el uso de liquidsoap, y hay dos fuentes (generada por IDJC) y 2 puntos de montaje con baja (64 kbps) y alta (128 kbps) calidad de sonido (configurando el bitrate de las fuentes en IDJC) en el formato ogg vorbis, Se recomienda hacer bien los cálculos para que se pueda dar la suficiente calidad si son pocos o muchos los usuarios.\cite{ref5} La desventaja de éste es que claramente no hay un sistema de respaldo y en caso de que el IDJC falle, nuestro punto de montaje desaparecerá, se recomienda tener otros fuentes que deberían ser administrados a través de liquidsoap.

\begin{entradilla} %codigo para una entradilla
{\em {\color{introcolor}{``Aquel que sacrifica la libertad por seguridad, no merece ninguna de las dos''}} -- Benjamin Franklin}
\end{entradilla}

El ideal para todos estos tipos de sistemas se encuentra en la wiki de RadioŃú la dirección de dicho esquema se encuentra en \url{http://wiki.radiognu.org/como_funciona_una_radio} éste esquema es la proyección a futuro de Radio GLUD; pongo la imagen para que se pueda apreciar rápidamente.

\begin{center}
\resizebox{8cm}{!}{{\epsfbox{images/radioglud/radiognu1.eps}}}
\mycaption{Esquema de como debería ser el flujo de datos en un sistema de streaming, extraído de \url{http://wiki.radiognu.org/como_funciona_una_radio}}
\end{center}

Si desean saber más de las tecnologías usadas en Radio GLUD, pueden visitar la dirección \url{http://radio/InformeFinaldeActividades-RadioGLUD1.odt} donde se explica detalladamente todo lo que se ha hecho y donde se puede descargar.

Para finalizar esta sección quiero dar crédito a los proyecto RadioŃú y Radio Liberación que ha dado mucho a la comunidad hispano-hablante ya que que con su documentación han dado un paso increíble para los que no sabemos leer muy bien Inglés (o definitivamente no sabemos).



% Siguiente página
\ebOpage{introcolor}{0.35}{Tecnologías libres}

\sectiontext{white}{black}{VII. Instrucciones para la primera emisión en una red local}

Ahora que ha terminado toda la teoría se hace necesario ir a lo técnico, pero antes hay que aclarar la notación que se va a utilizar. La distribución en que se hizo fue Fedora GNU/Linux 16 Verne con gnome 3, para evitar todos estos pasos se podría instalar una distribución como Musix GNU+Linux que es una distro dedicada a la música; la notación \verb!$ comando! representa los comandos en una terminal como usuario normal (sin permisos de root o administrador); la notación \verb!# comando! representa la ejecución de los comandos especificados como root, para ello solo basta con digitar \verb!$ su! y escribir la clave de root, en caso de que se maneje con la utilidad \verb!sudo! habría que digitarse \verb!$ sudo su! y la clave del usuario común; en particular debido a que \verb!sudo! es una herramienta muy conveniente para copiar y pegar comandos en terminal sin cambiar de usuario, es la que se utiliza en este mini-tutorial, pero si no se tiene, pueden loquearse como root y ejecutar solamente los comandos que están con sudo.

También es importante decir que los pasos no están tan detallados como se quisiera, pero se espera que con la teoría se pueda deducir lo que se dice que se haga o buscar en Internet para que quede claro, por ejemplo como se modifican los archivos con el editor nano, el lenguaje usado en los archivos o incluso como se abre una terminal y que es, ya que lastimosamente no hace parte de lo que se quiere en este artículo, pero por suerte hay mucha información en la red sobre estos temas.

%%%%%%%%%%FALTAN LOS PASOS%%%%%%%%%%%%%%

%%%%%%%%%%%%%%%%%%%%%%%%%%%%%%%%%%%%%%%%%%%%%%%%%%%%%%%%%%%%%%%%%%%
%\ebOpage{introcolor}{0.35}{Seccion dejar vacio} % Salto de página 1
%%%%%%%%%%%%%%%%%%%%%%%%%%%%%%%%%%%%%%%%%%%%%%%%%%%%%%%%%%%%%%%%%%%

\bibliographystyle{abbrv}
\begin{bibliografia}
\bibitem{ref1}
Copia textual del párrafo 3 del artículo Demonio (informática). Wikipedia la enciclopedia libre.
Recuperado de \url{es.wikipedia.org/wiki/Demonio_(informática)} el 19 de agosto de 2012.
\bibitem{ref2}
Artículo Música Libre. Wikipedia la enciclopedia libre.
Recuperado de \url{es.wikipedia.org/wiki/Música_libre} el 10 de agosto de 2012.
\bibitem{ref3}
Artículo Mickey Mouse. Wikipedia la enciclopedia libre.
Recuperado de \url{es.wikipedia.org/wiki/Mickey_Mouse} el 19 de agosto de 2012.
\bibitem{ref4}
Artículo Copyright Term Extension Act. Wikipedia the free encyclopedia.
Recuperado de \url{en.wikipedia.org/wiki/Sonny_Bono_Copyright_Term_Extension_Act} el 19 de agosto de 2012.
\bibitem{ref4}
Artículo Estándares para Calidad y Ancho de Banda. RadioŃú.
Recuperado de \url{wiki.radiognu.org/calidad} el 22 de agosto de 2012.
\end{bibliografia}


\begin{biografia}{images/radioglud/autor.eps}{Jorge Ulises Useche Cuellar} Estudiante de Ingeniería Electrónica, Interesado en el desarrollo de tecnologías libres que reduzcan la brecha social que existe hoy en día en Colombia al ser un país subdesarrollado. Viniendo de otro sector de Colombia encontró una simpatía inmediata al conocer el software libre e invierte su tiempo conociendo y desarrollando para la comunidad de software libre que le ha dado tanto.
\end{biografia}

\raggedcolumns
\pagebreak


\end{multicols}

\clearpage
\pagebreak

% Este fichero es parte del Número 5 de la Revista Occam's Razor
% Revista Occam's Razor Número 5
%
% (c)  2010 The Occam's Razor Team
%
% Esta obra está bajo una licencia Reconocimiento 2.5 España de Creative
% Commons. Para ver una copia de esta licencia, visite 
% http://creativecommons.org/licenses/by/2.5/es/
% o envie una carta a Creative Commons, 171 Second Street, Suite 300, 
% San Francisco, California 94105, USA.


\rput(2.5,-2.3){\resizebox{!}{3.2cm}{{\includegraphics[]{images/Importancia/white}}}}


%Imagen de el comienzo de el articulo, coordenadas desde %la parte superior izquierda del margen de la pagina

% -------------------------------------------------
% Cabecera
\begin{flushright}
\msection{introcolor}{black}{0.35}{Seguridad Informática}

\mtitle{10cm}{Importancia del Software Libre en la Seguridad Informática}

\msubtitle{8cm}{Importance of Free Software in Computer Security.}

{\sf Diego Osorio, Andrés Acosta, Leidy Aldana. \\Estudiantes de pregrado en Ingeniería de Sistemas}

{\psset{linecolor=black,linestyle=dotted}\psline(-12,0)}
\end{flushright}

\vspace{2mm}
% -------------------------------------------------

\begin{multicols}{2}

\sectiontext{white}{black}{Resumen}


% Introducción
\intro{introcolor}{E}{l software libre hace parte de una construcción social que beneficia a la comunidad; en el campo de la Seguridad Informática juega un papel importante con entornos y herramientas potentes que han revolucionado este campo. 
\\ 
%En esta parte por favor ubicar el Resumen de su artículo, debe permitir identificar en forma rápida el contenido básico del trabajo; no debe ser mayor a 250 palabras. No debe aportar información o conclusión que no está presente en el texto, así como tampoco debe  citar referencias bibliográficas. Debe quedar claro el problema que se investiga y el objetivo del mismo. 
}

\sectiontext{white}{black}{Palabras clave: } Confiabilidad, Integridad, Disponibilidad, Libertad, Seguridad.\\
\vspace{0.3cm}
% Leidy

% Introducción, Importancia del Software Libre, Seguridad Informática
\sectiontext{white}{black}{Introducción}

Compartir conocimiento ha aportado en la evolución de las sociedades, de forma similar, compartir código ha aportado a la construcción de software de mayor calidad; más allá del código abierto este artículo hace referencia al Software Libre, el cual permite la libertad del usuario y la libertad del software, permite modificar de acuerdo a necesidades propias, permite ayudar a una persona o a una comunidad; el software es libre al tener cuatro libertades, enunciadas en \textit{Software Libre para una Sociedad Libre} (Stallman, 2004, página 236): 

\vspace{0.1cm}

\begin{enumerate}
	\item[0.] La libertad cero es la libertad de ejecutar el programa con cualquier propósito. 
	\item[1.] La libertad uno es la libertad de ayudarte a ti mismo cambiando el programa para que se ajuste a tus necesidades. 
    \item[2.] La libertad dos es la libertad de ayudar al prójimo distribuyendo copias del programa. 
    \item[3.] La libertad tres es la de ayudar a distribuir a tu comunidad publicando una versión mejorada de modo que los otros puedan beneficiarse de tu trabajo.
\end{enumerate}

\vspace{0.15cm}

Gracias a estas cuatro libertades el Software Libre ha contribuido y puede contribuir en el aprendizaje, en la construcción y el soporte de grandes empresas o en la enseñanza; de hecho puede ser una alternativa que resulta más interesante, alejada del tedio y más próxima al área científica. Cabe resaltar que libre es diferente de gratis (libre como pensamiento libre más no como barra libre), que el software libre no quita el reconocimiento de la persona que contribuye a crear software o otro aporte hecho a la construcción de conocimiento libre.

\vspace{0.3cm}

\begin{entradilla}
{\em El {\color{introcolor}{Software Libre}} tiene la ventaja de ser estudiado, ejecutado, modificado y distribuido.}
\end{entradilla}

\vspace{0.3cm}

Su importancia se ve reflejada en la mejora continua, en que brinda la oportunidad de aprender apartir de aportes realizados por otras personas, se ve reflejada en el hecho de que el conocimiento debería ser libre. Para algunas personas más allá de estas ventajas mencionadas anteriormente, la elección de utilizar Software Libre es una cuestión ética, de ciencia, de no permanecer tanto en la ignorancia, tomando como ejemplo el área de la seguridad informática (área que hace parte de la mayoría de las tecnologías de la información y la comunicación), hay programas no libres que cuentan con código oculto utilizado para conseguir información del usuario o ejecutar procesos sin permiso del usuario.\\

Se denomina por seguridad informática a las medidas y controles que garantizan la confiabilidad, integridad y disponibilidad de los activos del sistema de información; incluyendo hardware y software, firmware e información procesada, almacenada y comunicada. \\

El área de la seguridad informática influye en la mayoría de las áreas en la actualidad, como en el desarrollo web, en la inteligencia artificial, en el manejo de grandes volúmenes de información; desde una perspectiva más cotidiana influye en el simple hecho de enviar un correo electrónico (para que llegue sin modificaciones exclusivamente al destinatario), de hacer una transacción bancaria en linea, o realizar la conexión a una red wi-fi sin que la información de nuestros dispositivos sea vulnerable.

% Siguiente página
%%%%%%%%%%%%%%%%%%%%%%%%%%%%%%%%%%%%%%%%%%%%%%%%%%%%%%%%%%%%
\ebOpage{introcolor}{0.35}{Seguridad informática}
%%%%%%%%%%%%%%%%%%%%%%%%%%%%%%%%%%%%%%%%%%%%%%%%%%%%%%%%%%%%

\begin{entradilla}
{\em La {\color{introcolor}{seguridad informática}} son las medidas y controles que garantizan la confiabilidad, integridad y disponibilidad de los activos del sistema de información.}
\end{entradilla}

De esta forma, este artículo busca mostrar al lector la importancia del trabajo conjunto de estos dos campos a través de ventajas, de herramientas, de alternativas.


% Diego -------------------------------------
\sectiontext{white}{black}{Herramientas}
% Implementación de herramientas
% confiabilidad, contribución, acceso, guias y documentación

Las herramientas utilizadas en la seguridad informática que son de código abierto, permiten que la comunidad examine el código, lo entiendan e incluso lo modifiquen, lo que permite a los usuarios aprender acerca del software, ajustarlo a sus necesidades e incluso contribuir. Esto lo hace más versátil e incluso agrega la confiabilidad de usarlo sabiendo que muchas personas han examinado el código y nos da la certeza que hace lo que debe, que en este no se encuentra código malicioso que pueda afectar a quien lo use y que ya ha sido probado.
\\
\\
El fácil acceso a todo tipo de usuarios, dónde no se tiene que pagar por utilizar, y el gran aumento del uso de las tecnologías de la información y el mundo digital, permite que la comunidad que utiliza software libre en la seguridad informática sea cada vez más, en los últimos años ha ganado mucha popularidad, incluso la demanda laboral ha aumentado drásticamente en los últimos años al igual que el uso de software libre a nivel profesional en la seguridad informática y personas especializadas en él.
\\
\\
Al ser probado constantemente por diferentes personas, con diferentes enfoques hace que se encuentren mayor número de problemas y fallos en menor tiempo que un software que sólo es mantenido por una empresa. Al haber más personas que pueden contribuir a mejorarlo, se ve una evolución constante haciendo las herramientas de software libre actualizadas, versátiles y altamente competitivas.
\\
\\
Es bastante la información que se encuentra en Internet de las herramientas de software libre en el campo de la seguridad informática, en los principales blogs se discute y comparte especialmente sobre software de código abierto. Tanto organizaciones como personas individuales comparten constantemente documentación, guías, tutoriales, vídeos, experiencias, preguntas, trucos, etc. generando una gran cantidad de material en múltiples idiomas permitiendo que casi cualquier persona con iniciativa y con algunos conocimientos básicos de la informática pueda aprender a utilizar estas herramientas.\\

%Andres

% Siguiente página
%%%%%%%%%%%%%%%%%%%%%%%%%%%%%%%%%%%%%%%%%%%%%%%%%%%%%%%%%%%%
%\ebOpage{introcolor}{0.35}{Seccion dejar vacio}
%%%%%%%%%%%%%%%%%%%%%%%%%%%%%%%%%%%%%%%%%%%%%%%%%%%%%%%%%%%%
%Principales herramientas de seguridad informática
	%Distribución GNU/LINUX de seguridad informática
    	%Herramientas de software libre.
\sectiontext{white}{black}{Kali Linux}

Es una distribución de software libre la cual esta en constante evolución con nuevas características que se agregan constantemente, esta distribución esta diseñada específicamente para pruebas de penetración profesional y auditoria de seguridad.\\

\sectiontext{white}{black}{Cambios principales de esta distribución} 

\sectiontext{white}{black}{Usuario único, acceso root por diseño:}

Es diseñado para trabajar en un único escenario root (usuario principal o administrador del sistema) ya que la mayoría de las herramientas utilizadas en pruebas de penetración requieren privilegios escalonados.\\

\sectiontext{white}{black}{Servicios de red deshabilitados de forma predeterminada:} 

Kali Linux contiene Demonios (scripts en segundo plano) que desactivan los servicios de red de forma predeterminada. Estos Demonios nos permiten instalar varios servicios en Kali Linux, a la vez que nos aseguramos de que nuestra distribución permanezca segura de forma predeterminada, sin importar qué paquetes estén instalados. Los servicios adicionales, como Bluetooth, también están en la lista negra de forma predeterminada.\\

\sectiontext{white}{black}{Núcleo Linux personalizado:} 

Kali Linux utiliza un kernel en sentido ascendente, con parches para inyección inalámbrica.\\

\sectiontext{white}{black}{Un conjunto mínimo y confiable de repositorios:} 

Dado los objetivos y metas de Kali Linux para mantener la integridad del sistema como un todo. Con ese objetivo en mente, el conjunto de fuentes de software ascendente que Kali utiliza se mantiene en un mínimo absoluto. Muchos nuevos usuarios de Kali están tentados a agregar repositorios adicionales a sus \textit{sources.list}, pero al hacerlo corre un riesgo muy serio de romper la instalación de Kali Linux.\\

\sectiontext{white}{black}{¿Es Kali Linux adecuado para usted?}

Los desarrolladores de Kali linux recomiendan que esta distribución debe ser usada por probadores de penetración profesional y especialistas en seguridad, ya que no es una distribución para soluciones generales

% Siguiente página
%%%%%%%%%%%%%%%%%%%%%%%%%%%%%%%%%%%%%%%%%%%%%%%%%%%%%%%%%%%%
\ebOpage{introcolor}{0.35}{Seguridad informática}
%%%%%%%%%%%%%%%%%%%%%%%%%%%%%%%%%%%%%%%%%%%%%%%%%%%%%%%%%%%%

de GNU/LINUX como el desarrollo de aplicaciones, diseño web, juegos, etc.\\

El uso indebido de herramientas de seguridad y de pruebas de penetración dentro de una red, en particular sin una autorización específica, puede causar daños irreparables y tener consecuencias significativas, personales y/o legales. \textit{"No entender lo que estabas haciendo"} no va a funcionar como una excusa.\\

Sin embargo, si usted es un probador de penetración profesional o está estudiando pruebas de penetración con el objetivo de convertirse en un profesional certificado, no hay mejor juego de herramientas - a cualquier precio - que Kali Linux.\\

\sectiontext{white}{black}{Algunas Herramientas de seguridad y hacking.} 

\sectiontext{white}{black}{Nmap \textit{("Mapeador de redes")}} 

Es una herramienta de código abierto para la exploración de redes y sondeo de seguridad de puertos, se diseño para analizar rápidamente grandes redes, utiliza paquetes IP "crudos" para poder determinar que equipos se encuentran disponibles en la red, esta herramienta generalmente se usa para auditoria de seguridad, administradores de redes, y es útil para realizar tareas rutinarias, como puede ser el inventariado de la red, la planificación de actualización de servicios y la motorización del tiempo que los equipos o servicios se mantiene activos en la red. El código fuente lo puedes encontrar en: \href{https://github.com/nmap/nmap}{github.com/nmap/nmap}\\

\sectiontext{white}{black}{John the Ripper} 

Es un programa libre de código abierto de criptografia, el cual aplica fuerza bruta para descifrar contraseñas, esta herramienta nos permite romper algoritmos de cifrado o hash, como DES, SHA-1 y otros. El código fuente lo puedes encontrar en: \href{https://github.com/magnumripper/JohnTheRipper}{github.com/magnumripper/JohnTheRipper}

\sectiontext{white}{black}{Nikto} 

Nikto es un software de código abierto (GPL), diseñado para evaluar la seguridad de los servidores web. Este software encuentra varios archivos predeterminados e inseguros, configuraciones y programas en cualquier tipo de servidor web. 
El código fuente lo puedes encontrar en: \href{https://github.com/sullo/nikto}{github.com/sullo/nikto}

\sectiontext{white}{black}{Wireshark} 

Wireshark es un software de código abierto (GPL) el cual lo usan administradores de redes para analizar el trafico de la red, capturar paquetes de diferentes tipos de protocolos que existen. El código fuente lo puedes encontrar en: \href{https://github.com/wireshark/wireshark}{github.com/wireshark/wireshark}

\sectiontext{white}{black}{Metasploit} 

Metasploit es un proyecto de código abierto para seguridad informática el cual proporciona información acerca de vulnerabilidades de seguridad y nos ayuda a realizar tests de penetración "Pentesting" junto con el desarrollo de firmas para sistemas de detección de intrusos. El código fuente lo puedes encontrar en:  \href{https://github.com/rapid7/metasploit-framework}{github.com/rapid7/metasploit-framework}

% Dudas respecto al manejo del artículo
% https://github.com/GLUD/articulo-individual-2017/blob/master/nombre_articulo.tex

\medskip

\bibliographystyle{abbrv}
\begin{bibliografia}
\bibitem{kopka}
Stallman, Richard. \emph{Software Libre para una Sociedad Libre}, \href{https://www.gnu.org/philosophy/fsfs/free_software.es.pdf}{disponible en pdf}.\hskip 1em plus
  0.5em minus 0.4em\relax Madrid, España: Traficantes de sueños, 2004.
  
  \bibitem{kopka}
2017 Offensive Security. \emph{Documentación oficial de Kali Linux}, \href{https://docs.kali.org/}{https://docs.kali.org/}.\hskip 1em plus
  0.5em minus 0.4em
  
\bibitem{kopka}
\emph{Guía de referencia de Nmap}, \href{https://nmap.org/man/es/}{https://nmap.org/man/es/}.\hskip 1em plus
  0.5em minus 0.4em
  
   \bibitem{kopka}
RAPID1. 
\emph{John the Ripper password cracker}, \href{http://www.openwall.com/john/}{http://www.openwall.com/john/}.\hskip 1em plus
  0.5em minus 0.4em
  
     \bibitem{kopka}
Wireshark Foundation . 
\emph{Documentación oficial de Wireshark}, \href{https://www.wireshark.org/docs/}{https://www.wireshark.org/docs/
}.\hskip 1em plus
  0.5em minus 0.4em
  

\end{bibliografia}


\begin{biografia}{images/Importancia/white.eps}{Diego Osorio} 
% añadir fotografía tamaño [2.5 cm x 3.3 cm ] 
%  |¬_¬|'  La foto es Opcional! , se añade reemplazando:
% \begin{biografia}{images/mi_articulo/autor.eps}{Nombre Completo del Autor}
Estudiante de Ingeniería de Sistemas, actual director del GLUD. Interesado en el Software Libre, el desarrollo de Software y la Seguridad Informática.\\daosoriog@gmail.com
\end{biografia}

\begin{biografia}{images/Importancia/white.eps}{Andrés Acosta} 
% añadir fotografía tamaño [2.5 cm x 3.3 cm ] 
%  |¬_¬|'  La foto es Opcional! , se añade reemplazando:
% \begin{biografia}{images/mi_articulo/autor.eps}{Nombre Completo del Autor}
Estudiante de Ingeniería de Sistemas.\\andres.taamap@gmail.com
\end{biografia}

\begin{biografia}{images/Importancia/white.eps}{Leidy Aldana} 
\\Estudiante de Ingeniería de Sistemas, interesada por aprender continuamente; quien ve en el software libre una perspectiva más agradable e interesante de su carrera.\\LeidyMarcelaAldana@gmail.com

\end{biografia}

\raggedcolumns
\pagebreak


\end{multicols}

\clearpage
\pagebreak

%% Este fichero es parte del Número 5 de la Revista Occam's Razor
% Revista Occam's Razor Número 5
%
% (c)  2010 The Occam's Razor Team
%
% Esta obra está bajo una licencia Reconocimiento 2.5 España de Creative
% Commons. Para ver una copia de esta licencia, visite 
% http://creativecommons.org/licenses/by/2.5/es/
% o envie una carta a Creative Commons, 171 Second Street, Suite 300, 
% San Francisco, California 94105, USA.

% Este .tex contiene el contenido del articulo
% Seccion (Dejar en blanco)
%

% Incluye imagen del artículo (Debe ser diferente del Kokopelli! (°_°) Sin embargo, no es obligatoria esta imagen)

%\rput(2.5,-2.3){\resizebox{!}{3.2cm}{{\includegraphics[]{images/nombre_articulo/glud}}}}


%Imagen de el comienzo de el articulo, coordenadas desde %la parte superior izquierda del margen de la pagina
% Las imagenes quedan mejor en formato .eps, agradecería si alguien logra colocar imagenes en dicho formato y comparte la corrección.

% -------------------------------------------------
% Cabecera
\begin{flushright}
\msection{introcolor}{black}{0.35}{Realidad Aumentada}

\mtitle{10cm}{TrascendentAR}

\msubtitle{8cm}{Librería de realidad aumentada para Android}

{\sf Juan Carlos Velandia, Estudiante de Ing. Electrónica\\
Marlon Arias Cárdenas, Estudiante de Ing. de Sistemas\\
Laydi Viviana Bautista, Estudiante de Ing. Catastral y Geodesia}

{\psset{linecolor=black,linestyle=dotted}\psline(-12,0)}
\end{flushright}

\vspace{2mm}
% -------------------------------------------------

\begin{multicols}{2}

\sectiontext{white}{black}{Resumen}

% Introducción
\intro{introcolor}{E}{n está oportunidad vamos a presentar una librería que integra fácilmente el framework de ARToolKit en el modulo de android de libGDX, Trascendent-AR, con el fin de ofrecer una herramienta de código abierto para el desarrollo de aplicaciones con Realidad Aumentada. Expondremos la metodología de trabajo que se desarrollo durante el proyecto, herramientas utilizadas como GitHub, Taiga.io y Jekyll, además de la documentación generada bajo el sitio web de trascendentAR. }

\sectiontext{white}{black}{Palabras clave: } Realidad Aumentada, libGDX, ARToolKit, Java, Android

% la introducción del texto. \\


\sectiontext{white}{black}{I. Introducción}\\
Las aplicaciones de realidad aumentada se han popularizado en los últimos años gracias al desarrollo tecnológico, esto lleva a que cada día más desarrolladores se interesen a incursionar en este campo, por lo que requieren herramientas especiales. En el mundo de los videojuegos los motores más populares han integrado librerías como Vuforia y ARToolKit que permitan a los usuarios existentes desarrollar Apps de realidad aumentada, dichos motores usualmente requieren trabajar en lenguajes como C\# o C++.\\ 

Con el objetivo de dar a los desarrolladores de Android una herramienta de código abierto que les permitiera crear poderosas aplicaciones de realidad aumentada de forma más simple, desde el grupo GNU/Linux de la Universidad Distrital nació Trascendent-AR, una librería que integra el framework de Realidad Aumentada ARToolKit con el framework de videojuegos libGDX.\\

En este punto, usted podría estar confundido con los conceptos de Framework, ARToolkit y libGDX, pero no se preocupe que se lo vamos a explicar, para empezar un framework podría traducirse aproximadamente como un “marco de trabajo” y es básicamente un esquema (un esqueleto, un patrón) para el desarrollo y/o la implementación de una aplicación.
\\

La \textbf {Realidad Aumentada} es la encargada de estudiar las técnicas que permiten integrar en tiempo real contenido digital con el mundo real, Milgram y Kishino la definen como los entornos en los que “se presentan objetos del mundo real y objetos virtuales de forma conjunta en una única pantalla”. \cite{libroRA}

%{Angulo}{ruta de la imagen}{escala}
\myfig{0}{images/TrascendentAR/simpleapp_output.png}{1}
\mycaption{Aplicación demo usando trascendentAR}

\textbf {ARToolKit} es una biblioteca de funciones para el desarrollo rápido de aplicaciones de Realidad Aumentada, entre las características destacables de ARToolKit se encuentra que es Multi-plataforma, es decir, que funciona en gran variedad de sistemas operativos, tiene una comunidad activa, y está bajo GNU General Public License (LGPLv3), a nivel técnico cabe resaltar que facilita el problema del registro de la cámara que es el proceso de posicionar la cámara (posición y rotación) relativamente a los objetos de la escena, dado que existen diferentes métodos orientados a diferentes tecnologías, ARToolKit emplea métodos de visión computacional, de forma que obtiene el posicionamiento relativo de 6 grados de libertad haciendo el seguimiento de marcadores cuadrados en tiempo real. 

\begin{entradilla}
{\em Herramienta de código abierto que les permitirá crear poderosas aplicaciones de {\color{introcolor}{Realidad Aumentada}} de forma más simple.}
\end{entradilla}
 
Continuando con los términos empleados inicialmente vamos a seguir con libGDX que se puede definir como  un framework para el desarrollo de videojuegos que provee un amplio abanico de plataformas finales, compatible con Windows, GNU/Linux, Android, iOS,  bajo licencia Apache2 escrito en lenguaje de programación Java. 
 
% Siguiente página
%%%%%%%%%%%%%%%%%%%%%%%%%%%%%%%%%%%%%%%%%%%%%%%%%%%%%%%%%%%%
\ebOpage{introcolor}{0.35}{Seccion dejar vacio}
%%%%%%%%%%%%%%%%%%%%%%%%%%%%%%%%%%%%%%%%%%%%%%%%%%%%%%%%%%%%

Es así como llegamos a TrascendentAR que ofrece una interfaz de fácil integración de ARToolKit en el módulo de Android de libGDX y se encuentra disponible para toda la comunidad de desarrolladores de videojuegos con libGDX bajo Licencia de Apache, Versión 2.0

% Fin de la introducción.\\

\sectiontext{white}{black}{II. Metodología}

Integrar ambos frameworks requiere inicialmente conocer los servicios que ofrece cada uno de ellos. Por ejemplo, ARToolKit provee entrenamiento y reconocimiento de marcadores, que son imágenes que pueden ser identificadas usando visión computacional, es decir, maneja todo el módulo de visión computacional y nos responde a preguntas como ¿Hay un marcador visible?, ¿Cuál marcador es?, ¿Qué coordenadas tiene dicho marcador en el espacio?
\\

Por otro lado libGDX provee de forma sencilla el manejo de gráficas, sonido, memoria, entradas táctiles, giroscopio, acelerómetro, vibración, y demás elementos que permiten crear aplicaciones interactivas en 2D y 3D.

\begin{entradilla}
{\em ... hacemos creer al ojo humano que {\color{introcolor}{existe}} un objeto en el mundo real.} 
\end{entradilla}

Después de este pequeño análisis se puede saber que la verdadera integración corresponde especialmente a la parte gráfica además del envío de señales por parte del módulo de realidad aumentada. 
\\

La parte complicada del proceso fue dibujar un objeto 3D en las coordenadas reales provistas por ARToolKit, antes de continuar les advertimos que esta sección va a ser un poco matemática ¿Recuerdan que ARToolKit ofrece las coordenadas reales del marcador? Bueno la forma en que se obtienen dichas coordenadas es en dos  matrices de 4x4, hay que mencionar que la programación de gráficos es pura y física álgebra lineal. 
\\

Usando dos matrices en conjunto podemos colocar un modelo 3D justo encima del marcador y mostrándolo correctamente en la pantalla hacemos creer al ojo humano que existe un objeto en el mundo real, tal como se ve en la figura 1. 
\\

La primera matriz es la de proyección de la cámara, transforma las coordenadas de los modelos a coordenadas 2D que corresponden a los pixeles de pantalla, o de forma inversa, tocando la pantalla (2D) proyectar las coordenadas sobre el espacio 3D de la aplicación permitiendo saber si se ha presionado sobre algún objeto. Una forma sencilla de entender la matriz de proyección, es dibujar un cubo sobre un papel; un cubo real es tridimensional pero su dibujo se encuentra en 2 dimensiones.
\\

\myfig{0}{images/TrascendentAR/projection_cube.png}{1}
\mycaption{Proyección de un cubo en 3D sobre un plano 2D}

\\

La segunda matriz se conoce como la matriz de trasformación del marcador, e indica la posición, rotación y escala del marcador en el espacio.

%\sectiontext{white}{black}{ Subsección de la Metodología} Especifica brevemente elementos y cómo se realiza, en caso de que pueda replicarse.

\sectiontext{white}{black}{Organización del proyecto}

TrascendentAR es un proyecto un tanto ambicioso, para culminarlo se debían cumplir un gran número de tareas que fueron asignadas a cada integrante del equipo, es así que es necesita conocer en que está trabajando tu compañero y que cambios ha hecho a la librería. Por esto utilizamos dos herramientas que facilitaron enormemente el desarrollo de TrascendentAR: GitHub y Taiga.io\\

\sectiontext{white}{black}{Github} es una plataforma de desarrollo colaborativo de software para alojar proyectos utilizando el sistema de control de versiones Git. Lo interesante de esta herramienta es que ofrece diversas funcionalidades en donde se encuentra una wiki para el mantenimiento de las distintas versiones de las páginas, un sistema de seguimiento de problemas que permiten a los miembros del equipo detallar un problema con el software, una herramienta de revisión de código, donde se pueden añadir anotaciones en cualquier punto de un fichero y debatir sobre determinados cambios realizados en un commit específico y un visor de ramas donde se pueden comparar los progresos realizados en las distintas ramas de nuestro repositorio. \cite{github}


\resizebox{8cm}{!}{{\myfig{0}{images/TrascendentAR/github-logo.png}{1}}}


% Siguiente página
%%%%%%%%%%%%%%%%%%%%%%%%%%%%%%%%%%%%%%%%%%%%%%%%%%%%%%%%%%%%
\ebOpage{introcolor}{0.35}{Seccion dejar vacio}
%%%%%%%%%%%%%%%%%%%%%%%%%%%%%%%%%%%%%%%%%%%%%%%%%%%%%%%%%%%%

\sectiontext{white}{black}{Taiga.io} es una aplicación que gestiona el progreso de un proyecto. Permite integrarse con repositorios de control de versiones basados en la web como GitHub y utiliza metodologías de desarrollo de software ágil e incremental para gestionar el desarrollo de productos como Scrum o Kanban. 

Kanban proveniente de Japón, y se refiere más a una metodología en donde se divide el desarrollo de proyectos en etapas. Es así como Taiga.io se convierte en una gran herramienta de apoyo ofreciendo un tablero Kanban que se utilizamos para gestionar las tareas asignadas a cada uno de los integrantes por medio de tarjetas y estados de avance del proyecto. 

\sectiontext{white}{black}{III. Distribución}

Uno de los objetivos de este proyecto es que sea utilizado por otros desarrolladores de aplicaciones, además de que sea relevante y trascendental... Es aquí en donde una vez finalizado el proyecto trabajamos en una excelente documentación y un sitio web atractivo. 


\sectiontext{white}{black}{Documentación}\\
TrascendentAR en si mismo comprende solamente tres clases de Java, por lo que la cantidad de líneas de código es pequeña, cerca de 400 líneas aunque trabajamos bajo la filosofía de que el código en si mismo debe ser autoexplicativo y fácil de entender tomamos la decisión de complementar la documentación por medio de tutoriales publicados en el sitio web del proyecto.

\myfig{0}{images/TrascendentAR/website.png}{1}
\mycaption{Sitio web de TrascendentAR}

\sectiontext{white}{black}{Sitio Web}\\
Con respecto al sitio web, sabemos que una buena primera impresión cuenta más de lo que imaginamos, por esa razón decidimos que la portada de TrascendentAR debía ser llamativa y el sitio web debía ser más que una aburrida página en blanco de Github, aún así, existían varios desafíos por considerar:

\begin{itemize}
    \item \textbf {Tiempo.} No debíamos excedernos más de dos semanas en tener un sitio web funcional y con contenido.
    \item \textbf {Originalidad.} Elegir un Nombre de Dominio confiable y profesional.
    \item \textbf {Visibilidad.} La página debía tener atractivo visual.
    \item \textbf {Accesibilidad.} Ofrecer una curva de aprendizaje suave.
    
\end{itemize}

La respuesta de la mayoría de los desafío la encontramos en el uso de Jekyll en conjunto con Github Pages.
\\

\textbf {Jekyll} es un generador de sitios web estáticos y blogs partir de componentes dinámicos como plantillas además de permitir convertir documentos en texto plano en páginas web cuidadosamente diseñadas. En contraste, \textbf {Github Pages} es un servicio de host provisto por Github y así, desde el mismo repositorio del proyecto en GitHub se genera el sitio web en jekyll, para el caso de Trascendentar el sitio web se despliega en la siguiente dirección  {\color{introcolor}{glud.github.io/trascendentAR}} la cual resulta ser para nuestros propósitos una dirección confiable y lo suficientemente profesional.
\\ 
\begin{entradilla}
{\em Un excelente producto en si mismo no vale nada... {\color{introcolor}{si nadie lo conoce}} }
\end{entradilla}
\\

Es necesario señalar que no creamos el sitio web desde cero, por el contrario adaptamos una plantilla de jekyll tomada de {\color{introcolor}{www.jekyllthemes.org}} así logramos ahorrar tiempo además de obtener un sitio web bien diseñado y fácilmente adaptable a diferentes dimensiones de pantalla en un corto periodo de tiempo. 
\\

No obstante si modificamos en gran medida el código fuente del sitio web, adaptando estilos y añadiendo soporte para otros idiomas obteniendo así mayor alcance en la comunidad, esto lo hemos evidenciado en nuestros primeros usuarios que no saben Español pero si Inglés. Actualmente el sitio web está disponible en Inglés y Español, en donde se pueden encontrar diferentes tutoriales para utilizar TranscendertAR, las especificaciones de la licencia y algunas opciones para contribuir al proyecto en donde también se ofrece el código fuente.



\bibliographystyle{abbrv}
\begin{bibliografia}
\bibitem{libroRA}
C. Gonzales, D. Vallejo, J. Albusaz \& J, Castro. \emph{Realidad Aumentada. Un Enfoque Práctico con ARToolKit y Blender}, 1rd~ed.\hskip 1em plus
  0.5em minus 0.4em\relax ISBN: 978-84-686-1151-8. 2012. 
  
\bibitem{github}
 L. Castillo. \emph{¿Qué es GitHub?}, \hskip [Blog Online]
 Recuperado de: http://conociendogithub.readthedocs.io/

\bibitem{apache2}
 L. Castillo. \emph{¿Qué es GitHub?}, \hskip [Blog Online]
 Recuperado de: http://conociendogithub.readthedocs.io/

\end{bibliografia}

%%%%%AUTORES%%%%%%%%%%%%%%%%%5

% añadir fotografía tamaño [2.5 cm x 3.3 cm ] 
%  |¬_¬|'  La foto es Opcional! , se añade reemplazando:
% \begin{biografia}{images/mi_articulo/autor.eps}{Nombre Completo del Autor} corta descripción \end{biografia}

\begin{biografia}{}{Juan Carlos Velandia} 
Estudiante de Ing. Electrónica.\\
\textbf {Laydi Viviana Bautista} Estudiante de Ing. Catastral y Geodesia.\\
\textbf {Marlon Arias Cárdena} Estudiante de Ing. de Sistemas.\\

Estudiantes de la Universidad Distrital FJC, Miembros Activos del GLUD - Grupo GNU/Linux Universidad Distrital
\end{biografia}




\raggedcolumns
\pagebreak


\end{multicols}

\clearpage
\pagebreak


% \begin{center}
% \resizebox{5cm}{!}{{\epsfbox{images/nombre_articulo/glud.eps}}}
% \mycaption{Gráfico del movimiento de los datos multimedia desde la generación del fuente hasta que lo escucha el usuario o oyente de la radio o canal de vídeo}
% \end{center}

%Para añadir un link:
%\href{https://glud.org/}{Hipervínculo} \\
%\href{https://blog.desdelinux.net/occamss-razor-publicacion-libre-de-divulgacion-cientifica/}{Acerca de Occam's razor} \\

%\textsf {\textbf {Código original }}


% Este fichero es parte del Número 5 de la Revista Occam's Razor
% Revista Occam's Razor Número 5
%
% (c)  2010 The Occam's Razor Team
%
% Esta obra está bajo una licencia Reconocimiento 2.5 España de Creative
% Commons. Para ver una copia de esta licencia, visite 
% http://creativecommons.org/licenses/by/2.5/es/
% o envie una carta a Creative Commons, 171 Second Street, Suite 300, 
% San Francisco, California 94105, USA.

% Este .tex contiene el contenido del articulo
% Seccion (Dejar en blanco)
%

% Incluye imagen del artículo (Debe ser diferente del Kokopelli! (°_°) Sin embargo, no es obligatoria esta imagen)

%\rput(2.5,-2.3){\resizebox{!}{3.2cm}{{\includegraphics[]{images/nombre_articulo/glud}}}}


%Imagen de el comienzo de el articulo, coordenadas desde %la parte superior izquierda del margen de la pagina
% Las imagenes quedan mejor en formato .eps, agradecería si alguien logra colocar imagenes en dicho formato y comparte la corrección.

% -------------------------------------------------
% Cabecera
\begin{flushright}
\msection{introcolor}{black}{0.35}{Accesibilidad de plataformas}

\mtitle{10cm}{El aporte del software libre a la accesibilidad en plataformas tecnológicas.}

\msubtitle{8cm}{The contribution of free software to accessibility in technology platforms.}

{\sf Jesús David Romero, Estudiante de Ingeniería de Sistemas }

{\psset{linecolor=black,linestyle=dotted}\psline(-12,0)}
\end{flushright}

\vspace{2mm}
% -------------------------------------------------

\begin{multicols}{2}

\sectiontext{white}{black}{Resumen}

% Introducción
\intro{introcolor}{E}{ste trabajo hace un estado de arte sobre los avances en el software libre, centrándose en las herramientas que contribuyen a la accesibilidad de las plataformas tecnológicas. 
}

\sectiontext{white}{black}{Palabras clave: } Software libre, accesibilidad, plataformas tecnológicas, avance, discapacidad.


\sectiontext{white}{black}{Introducción}

La accesibilidad es una tendencia que pretende brindarle a las personas que presenten algún tipo de discapacidad, ya sea visual, auditiva o cognitiva una experiencia parcialmente completa a los contenidos audiovisuales presentes en las diferentes plataformas tecnológicas, aliado a la ideología emanada por el software libre las oportunidades de avance que hay en este campo son tan limitadas como la imaginación de los colaboradores, y la oportunidad de generar cada vez mayor inmersión en estos contenidos creando una equidad en el mundo de las tecnologías.
 
\noindent

\sectiontext{white}{black}{Accesibilidad en las tecnologías de la información}

La accesibilidad se basa en el diseño de producto y entornos de fácil uso para el mayor número de personas, enfocandose a la población más vulnerable en cuanto al uso de herramientas se refiere. Existen siete principios que apuntan al tema de la accesibilidad: la igualdad de uso, la flexibilidad, diseño simple e intuitivo, información perceptible, diseño tolerante a errores, diseño con escaso esfuerzo físico y tamaño y espacio para el acceso y uso.

Específicamente la WGAC tiene tres principios: Debe ser perceptible, operable, robusto.

La accesibilidad recoge varias clases de discapacidades: visual, cognitiva, auditiva y neurológica. Recoge también a personas que no poseen ningún tipo de discapacidad pero que por alguna condición no pueden disfrutar por completo su experiencia en la web (por ejemplo, la conexión lenta, computadores de bajos recursos, pérdida momentánea de un brazo, edad avanzada).

La W3C define accesibilidad como: \textit{"Acceso de todos a la Web, independientemente del tipo de hardware, software, infraestructura de red, idioma, cultura, localización geográfica y capacidades de los usuarios"}.


\sectiontext{white}{black}{Software Libre}

\begin{quote}
\raggedleft
\textit{“ Enseñar a los niños el uso del software libre en las escuelas, formará individuos con sentido de libertad. "}\\
 Richard Stallman

\end{quote}

El software libre es el software que respeta la libertad de los usuarios y de la comunidad, en otras palabras es el derecho y la libertad que tienen los usuarios de ejecutar, copiar, distribuir, estudiar, modificar y mejorar el software; entendiendo así el software libre como una cuestión de libertad y no de precio ya que dada su escritura en ingles da para equivocaciones y malas interpretaciones.

Se dice que un software es libre cuando cumple con las siguientes 4 libertades:
\begin{itemize}
\item (Libertad 0) La libertad de ejecutar el programa como se desea, con cualquier propósito.
\item (Libertad 1) La libertad de estudiar cómo funciona el programa, y cambiarlo para que haga lo que usted quiera. El acceso al código fuente es una condición necesaria para ello.
\item (Libertad 2) La libertad de redistribuir copias para ayudar a su prójimo.
\item (Libertad 3)La libertad de distribuir copias de sus versiones modificadas a terceros. Esto le permite ofrecer a toda la comunidad la oportunidad de beneficiarse de las modificaciones. El acceso al código fuente es una condición necesaria para ello.
\end{itemize}

% Siguiente página
%%%%%%%%%%%%%%%%%%%%%%%%%%%%%%%%%%%%%%%%%%%%%%%%%%%%%%%%%%%%
\ebOpage{introcolor}{0.35}{Accesibilidad de plataformas}
%%%%%%%%%%%%%%%%%%%%%%%%%%%%%%%%%%%%%%%%%%%%%%%%%%%%%%%%%%%%


\sectiontext{white}{black}{Ventajas del Software libre}

Cuando se habla de libertad en el software libre y no se tiene profundidad en el tema, se puede quedar en el campo de la filosofía abstracta y poco práctica, por eso a continuación se esbozan algunas ventajas identificadas en el campo de la implementación de esta praxys filosófica y revolucionaria.

\begin{itemize}
\item Revisión pública: Al estar varias personas revisando sin ningún tipo de restricción el código, la detección de errores y bugs va a ser mucho más rápida.


\item Distribución de conocimientos: Aunque se es enfático en que software libre no significa software gratuito, si hay casos en que efectivamente este no presenta precio para adquirirlo, permitiendo que la adquisición de conocimientos sea mucho más dinamica y más nutrida.
\item Innovación: En el libro \textit{La catedral y el bazar} de Erick Raymond, se muestra la diferencia entre un modelo de construcción a modo catedral, hermético y vertical y un modelo bazar, variado horizontal y “bullicioso” haciendo analogía con el software privativo y el software libre respectivamente, este modelo bazar donde hay un caos constante y un ir y venir de ideas e iniciativas permite efectivamente que de todo este caos surgen constantemente mejoras y nuevos desarrollos al software inicial haciendo que la innovación esté al orden del dia.
\item Beneficio económico: Como ya se mencionó antes en efecto mucho software libre es totalmente gratuito, permitiendo que en empresas la inversión en informática no sea tan alta.
\item Seguridad: Esta ventaja va de la mano de la revisión pública pero se hace el enfásis ya que con la cantidad de datos que se manejan en la web la seguridad debe ser un indicador primordial.
\end{itemize}

\sectiontext{white}{black}{Proyectos de Software libre y accesibilidad}
\subsection*{Discapacidad visual}
\begin{itemize}
\item El proyecto Orca, para el sistema operativo Linux, combina herramientas de síntesis de voz para que el ordenador lea en voz alta lo que aparece en la pantalla, con la posibilidad de trabajar con Braille y magnificación de pantalla. Es parte de la plataforma “Gnome" de Linux. 
http://www.gnome.org/projets/orca
\item Lazaroux, este proyecto va enfocado a personas con discapacidad visual de habla hispana, es una distribución de Linux que ofrece varias herramientas y aplicaciones para facilitar la accesibilidad y un motor totalmente en castellano. No es necesaria su instalación para su utilización.
\item Mozilla Firefox, es un conocido proyecto de software libre que tiene una gran cantidad de extensiones para proporcionar accesibilidad en el contenido web.
\item Britty, es un proceso que se ejecuta en segundo plano y permite la entrada de teclados braille para usarlos en consolas de texto de Linux y Unix.
\item Festival, es un sintetizador de voces, disponible en inglés y en español para diferentes distribuciones. Con documentación completa para la implementación de nuevas voces.
\item Gnome-Speech, es una libreria API general que permite la programación de software basado en Gnome para producir voz, actualmente solo está disponible la interfaz de Festival.
\item Gnopernicus, brinda a los usuarios con discapacidad visual o visión limitada usar las aplicaciones en Gnome mediante herramientas como lupas magnificadoras, sintetizador de voz y lector de braille.
\item Screader, en un sintetizador de voz que reproduce texto y caracteres de la consola.
\item SVGATextMode, ajusta el tamaño de las líneas de texto en consola en tarjetas SVGA para Linux en modo texto. Modifica el tamaño de la fuente, el cursor, la sincronización de horizontal y vertical, etc. En modo texto se puede sacar todo el partido a la tarjeta de video y del monitor. 
\item KeyTouch, permite generar funciones extras en el teclado para ciertas operaciones, aunque el proyecto no va generado principalmente a personas con discapacidad, si se muestra un gran interés por esta población.
\end{itemize}

\subsection*{Discapacidad Motriz}
\begin{itemize}
\item Dasher, es una interfaz de texto que permite mediante acciones predictivas del ratón reemplazar la escritura mediante el teclado; utiliza inteligencia artificial y estadísticas de grupos de palabras e idiomas.
\item XVoice, permite el control total de las aplicaciones mediante la voz utilizando IBM's ViaVoice for Linux, permite desde controlar programas hasta comandos de usuarios.
\end{itemize}

% Siguiente página
%%%%%%%%%%%%%%%%%%%%%%%%%%%%%%%%%%%%%%%%%%%%%%%%%%%%%%%%%%%%
\ebOpage{introcolor}{0.35}{Accesibilidad de plataformas}
%%%%%%%%%%%%%%%%%%%%%%%%%%%%%%%%%%%%%%%%%%%%%%%%%%%%%%%%%%%%


\subsection*{Herramientas de educación}
\begin{itemize}
\item Clic, es un conocido software educativo que permite realizar diversas actividades: puzzles, asociaciones, crucigramas, sopas de letras, etc. Posee varias oportunidades de personalización. Jclic es la última versión, desarrollado como sw libre y disponible para MAC OS, Solaris, Windows y Linux, algunas características de click no estaban presentes en jclick pero gracias a su naturaleza de sw libre, varios grupos de colaboradores ya están incluyendolas.
\end{itemize}

\sectiontext{white}{black}{Conclusiones}\\
La accesibilidad es una tendencia que busca la equidad y el disfrute igualitario de los usuarios frente al contenido informatico, dado su auge se vuelve necesario (no por capricho) la elaboración de múltiples herramientas que respondan a estas necesidades, aparte de la necesidad de su desarrollo se requiere la democratización de este conocimiento ya que es un derecho, el derecho a informarse y la falta de estrategias para difundir este conocimiento no puede ser la excusa. Es aquí donde entra el sw libre permitiendo mediante el libre y constante discernimiento la creación de nuevas y más y mejores herramientas que proporcionen al usuario final ese derecho al disfrute total de la información. Esto lo demuestra el hecho de que a pesar de que a esta iniciativa se sumó el sw libre un poco tarde ya ha hecho avances bastante significativos, que hacen pensar que en un futuro no muy lejano llegará a superar a los sw privativos, tanto en funcionalidad como en etica.

\medskip

\bibliographystyle{abbrv}
\begin{bibliografia}
\textbf{Bibliografía}\\

Pautas, métodos y herramientas de evaluación de accesibilidad web Autor: Cinthia De Oleo Moreta, Luis Rodriguéz Baena. Revista Universidad de Manizales. 21 de junio del 2013.

\end{bibliografia}

\textbf{Cibergrafía}

\begin{enumerate}

\item Software libre y accesibilidad: Matías Sánchez Caballero. 15 de julio del 2010. Tomado de \href{http://www.nosolousabilidad.com/articulos/software_libre.htm?utm_source=feedburner}{www.nosolousabilidad.com}
% \href{https://www.wikibooks.org}{Wikibooks home}

\item El sistema operativo GNU, ¿Qué es el Software Libre? \href{https://www.gnu.org/philosophy/free-sw.es.html}{www.gnu.org}

\item Qué es la accesibilidad web \href{http://www.w3c.es/Traducciones/es/WAI/intro/accessibility}{www.w3c.es}

\item Mozilla Firefox, accesibilidad web \href{http://www.mozilla.org/access}{www.mozilla.org}

\item Festival \href{http://www.cstr.ed.ac.uk/projects/festival/}{www.cstr.ed.ac.uk}

\item GnomeSpeech \href{http://www.escomposlinux.org/lfs-es/blfs-es-SVN/gnome/gnome-speech.html}{www.escomposlinux.org}

\item Gnopernicus \href{http://www.escomposlinux.org/lfs-es/blfs-es-6.0/gnome/gnopernicus.html}{www.escomposlinux.org}

\item Screader \href{http://leb.net/pub/blinux/screader/}{leb.net}

\item SVGATextMode \href{http://freshmeat.net/projects/svgatextmode/}{freshmeat.net}

\item KeyTouch \href{http://keytouch.sourceforge.net}{keytouch.sourceforge.net}

\item Dasher \href{http://www.inference.phy.cam.ac.uk/dasher}{www.inference.phy.cam.ac.uk}

\item XVoyce \href{http://xvoice.sourceforge.net/}{xvoice.sourceforge.net}

\item Proyecto Clic \href{http://clic.xtec.es/es/jclic}{clic.xtec.es}

\item Britty \href{http://mielke.cc/brltty/index.html}{mielke.cc}

\item Proyecto orca \href{http://www.gnome.org/projets/orca}{www.gnome.org}



\end{enumerate}

\begin{biografia}{images/accesibilidad/white.eps}{Jesús David Romero} 
% añadir fotografía tamaño [2.5 cm x 3.3 cm ] 
%  |¬_¬|'  La foto es Opcional! , se añade reemplazando:
% \begin{biografia}{images/mi_articulo/autor.eps}{Nombre Completo del Autor}
\\
Estudiante de Ingeniería de Sistemas, miembro activo del Grupo GNU Linux Universidad Distrital
\end{biografia}

\raggedcolumns
\pagebreak


\end{multicols}

\clearpage
\pagebreak


\end{document}
